% concatenated from Jacobian_Algebras_of_Species_with_Potentials_and_2-Representation_Finite.tex
% !TeX spellcheck = en_GB
\documentclass[a4paper, reqno]{amsart}
\usepackage[utf8x]{inputenc} 
\usepackage{amsmath}
\usepackage{mathtools}
\usepackage{amssymb}
\usepackage{subcaption}
\usepackage{hyperref} % [hidelinks]
\usepackage{bookmark}
\usepackage{amsthm}
\usepackage{yhmath}
\usepackage{float}
\usepackage[absolute,overlay]{textpos}
\usepackage{titletoc}
%\usepackage{natbib}
\usepackage{tikz-cd}
\usepackage[left=2cm, right=2cm]{geometry}
\usepackage{enumitem}
\usepackage{preview}
\usetikzlibrary{bbox}
\usepackage{adjustbox}

\usepackage{xcolor}

\def\Image#1{\raisebox{-.5\height}{\includegraphics[scale=0.7]{#1}}}

\tikzcdset{
  cells={font=\everymath\expandafter{\the\everymath\displaystyle}},
}

\newtheorem{mythm}{Theorem}[section]
\newtheorem*{mythm*}{Theorem}
\newtheorem{myprop}[mythm]{Proposition}
\newtheorem*{myprop*}{Proposition}
\newtheorem{mylemma}[mythm]{Lemma}
\newtheorem{mycor}[mythm]{Corollary}
\newtheorem*{mycor*}{Corollary}

\theoremstyle{definition}
\newtheorem{mydef}[mythm]{Definition}
\newtheorem{myrem}[mythm]{Remark}
\newtheorem{myex}[mythm]{Example}
\newtheorem{myconjecture}[mythm]{Conjecture}
\newtheorem{mydefprop}[mythm]{Definition/Proposition}
\newtheorem{myconst}[mythm]{Construction}
\newtheorem{myquestion}[mythm]{Question}

\renewcommand{\theequation}{\arabic{section}.\arabic{mythm}.\arabic{equation}}
\counterwithin*{equation}{mythm}

\newcommand\norm[2]{C^{#2}\left(#1\right)}
\newcommand\note[1]{\adjustbox{bgcolor=yellow,minipage=[t]{\linewidth}}{\textbf{\Large #1}}}
\newcommand\lsub[1]{{}_{#1}\hspace{-0.25mm}}
\newcommand{\q}{\mathbb{Q}[\eta]}
\newcommand{\C}{\mathbb{C}}
\newcommand{\Z}{\mathbb{Z}}
\newcommand{\Q}{\mathbb{Q}}
\newcommand{\N}{\mathbb{N}}
\newcommand{\K}{\mathbb{K}}
\newcommand{\rk}{\mathrm{rk}\,}
\newcommand{\id}{\mathrm{Id}}
\newcommand{\R}{\mathbb{R}}
\newcommand{\chr}{\mathrm{char}\,}
\newcommand{\cO}{\mathcal{O}}
\renewcommand{\P}{\mathbb{P}}
\newcommand{\bd}{\partial}
\renewcommand{\t}{\mathfrak{t}}
\renewcommand{\b}{\mathfrak{b}}
\newcommand{\U}{\mathcal{U}}
\newcommand{\D}{\mathcal{D}}
\newcommand{\add}{\mathrm{add}}
\newcommand{\ind}{\mathrm{ind}}
\newcommand{\rad}{\mathrm{rad}}

\newcommand{\bigslant}[2]{{\raisebox{.2em}{$#1$}\left/\raisebox{-.2em}{$#2$}\right.}}

\DeclareMathOperator{\coker}{\text{coker}}
\DeclareMathOperator{\im}{\text{im}}
\DeclareMathOperator{\Hom}{\mathrm{Hom}}
\DeclareMathOperator{\End}{\mathrm{End}}
\DeclareMathOperator{\gen}{\text{Gen}}
\DeclareMathOperator{\gr}{\text{gr}}
\DeclareMathOperator{\sgn}{\mathrm{sgn}}
\DeclareMathOperator{\Tr}{\mathrm{Tr}}
\DeclareMathOperator{\Mat}{\mathrm{Mat}}
\DeclareMathOperator{\gldim}{\mathrm{gldim}}
\DeclareMathOperator{\type}{\mathrm{Type}}
\DeclareMathOperator{\supp}{\mathrm{supp}}
\DeclareMathOperator{\DD}{\mathsf{D}}

%\setlength\parindent{0pt}
%\setlength\parskip{0.5em}

%\setlength{\abovedisplayskip}{5pt}
%\setlength{\belowdisplayskip}{5000pt}
%\setlength{\abovedisplayshortskip}{2pt}
%\setlength{\belowdisplayshortskip}{2pt}

% \dottedcontents{<section>}[<left>]{<above-code>}
% {<label width>}{<leader width>}
%\dottedcontents{section}[1.3em]{\bfseries}{1em}{1pc}
%\dottedcontents{subsection}[0em]{}{3.3em}{1pc}

% center the toc heading
\renewcommand{\contentsname}{\centering Contents}
\title[Jacobian Algebras of Species with Potentials and $2$-Representation Finite Algebras]{Jacobian Algebras of Species with Potentials and $2$-Representation Finite Algebras}
\author{Christoffer S\"oderberg}

\begin{document}
\begin{abstract}
	We study $2$-representation finite $\K$-algebras obtained from tensor products of tensor algebras of species. In earlier work \cite{soderberg2024mutation} we computed the higher preprojective algebra of said algebras to be given as Jacobian algebras of certain species with potential $(S, W)$, which are self-injective and finite dimensional. Truncating these Jacobian algebras yields a rich source of $2$-representation finite $\K$-algebras. Under suitable assumptions, we prove that the set of all cuts of $(S, W)$ is transitive under successive cut-mutations. Furthermore, we show that cuts and cut-mutation correspond to truncated Jacobian algebras and $2$-APR tilting, respectively. Consequently, under certain assumptions, all truncated Jacobian algebras are related to each other via $2$-APR tilting. We produce various new examples of $2$-representation finite $\K$-algebras.
\end{abstract}

\maketitle
\tableofcontents

\iffalse
Things to do:
- Prove kellers recipe.
- mention bimodule structure of transpose.
- introduce |alpha| notation
- Fix Corollary 6.8 and Example 6.8.
- Make a plan for section 7 examples.
\fi

\iffalse
Structure of article
2. quiver with cycles and its canvas
- simply conn covered and enough cuts -> transitive
3. tensor products of quiver with cycles
- Q finite trees, Q x Q trans
- D x D algebra (S, W) trans
- F conn species, S x S morita trans
4. truncated jacobian algebras
- introduce alg cuts
- every cut is alg for selfinjective
5. 2 rep and truncated
- kellers recipe?
- alg cut -> 2 rep algebra
6. apr tils and cut-mutation
- apr tilt corresponds to cut-mutation
- l hom species tensor truncated 2 apr tilts from each other
	- change this since every representation finite l homogeneous species is F connected.

Random notes:
\begin{enumerate}
	\item[\cite{HerschendOsamu2011quiverwithpotential}] qps is a vital part of cluster algebras/categories. They are closely related to the $3$-Calabi Yau Ginzburg algebras. Taking the zeroth homology yields the jacobian algebra. Self-injective property is kept when mutating along Nakayama orbits. First main result asserts that $2$-rf algebras are given by Jacobian algebras of qps. Second main result says that self-inj preserves along mutationg along orbits. Third main results gives a sufficient condition for when qps are derived equivalent, saying that trunacted jacobian are $2$-apr from each other. Introduces canvas of qp.
	\item[\cite{dlab1975algebras}] Provides a list of representation finite species classified using the Dynkin diagrams.
	\item[\cite{Dlab_1980}] Defines Casimir elements for species and gives the construction of the preprojective algebra using said CAsimir elements.
	\item[\cite{Nquefack2012PotentialsJacobian}] Studies a more generalised version of Jacobian algebras. Being more explicit lets us define the Jacobian algebra for species with potentials.
	\item[\cite{IO11nRFalgandnAPRtilt}] Introduces the notion of $n$-apr tilting for $n$-representation finite hereditary algebras. They showed that $n$-apr tilting preserves $n$-representation finiteness.
	\item[\cite{IO13Stablecateofhigherpreproj}] Introduces the $(n+1)$-preprojective algebra of algebras of global dimension $n$. They show that if an algebra if $n$-representation finite then its $(n+1)$-preprojective algebra is self-injective. In this situation the stable module category of $(n+1)$-preprojective algebra is $(n+1)$-CY. More precisely, it is the $(n+1)$-Amiot cluster category of the stable $n$-Auslander algebra of the original algebra.
\end{enumerate}

\note{can we describe the other 2-apr tilts of tensor products, put in introduction and note it here}

around 6 lines about standard apr tilt (1 instead of 2) and mirror functors. cite dlab ringel where they talk about preprojective algebra. dont need to mention degree shift. Origin of preprojective algebra of quiver case is litele bit unclear. Maybe draw a diaram. 

build up towards quiver with potentials since it is not that common.

describe potential with distinguished cycles. 

interpret 2-apr as cut mutation (refer forward to corollary).

thn we introduce transitivity.

2-apr tilted algebras are strongly related, i.e. same preprojective algebra, derived equivalent etc. Thus a natural question is whether we have transitivity.

Our assumtpion that we have a cut and it described the preproecíve algebra and after cut-mutaiton we have the same preprojectie algebra. After transitivity we note that this assumption works for every cut if it works for one of them.

darius called it a preprojective cut, i.e. that the algebra is the algebra given by degree zero of preprojective algebra. In case degree 0 has gldim 2 and preprojectiev gien by jacobian, then it is called a preprojective cut.

If we use 2-apr tilt on A corresponds to mutation on cut and degree 0 part of preprojective is the 2-apr tilt. If we know that sp is such that it is transitive and we also know that all cuts are preprojective. 

The main focus of this article lies with the Jacobian algebras that is self-injective and finite dimensional. For a quiver with potential $(Q, W)$ one defines the Jacobian algebra as the path algebra of $Q$ modulo the ideal generated by derivations of $W$. This give rise to interesting geometric and algebraic properties, including cluster algebras [??] and categorification [??]. Jacobian algebras are equal to the zeroth homology of the Ginzburg algebra \cite{ginzburg2007calabiyaualgebras} associated with a quiver with potential, which is useful when studying properties of Calabi--Yau algebras. 

For a set of arrows $C$ in $Q$ we define a grading where we put all arrows in $C$ in degree $1$ and the rest of the arrows in degree $0$. This extends to a grading on $S$. We say that $C$ is a cut if the potential $W$ is homogeneous of degree $1$ in the grading induced by $C$. Given a cut $C$ we can define the truncated Jacobian algebra $\mathcal{P}(S, W, C)$ as the quotient of the Jacobian algebra $\mathcal{P}(S, W)$ with the ideal generated by $C$.

Truncated Jacobian algebras appears when studying $2$-representation finite $\K$-algebras. For the case of quivers with potential \cite{HerschendOsamu2011quiverwithpotential} showed that every basic $2$-representation finite $\K$-algebra of an algebraically closed field $\K$ is isomorphic to a truncated Jacobian algebra. Even with this, it is difficult to produce examples of $2$-representation finite $\K$-algebras in general. Our main theorem aims to produce an algoritm of how one can produce many new examples of $2$-representation finite $\K$-algebras.

\fi

\section{Introduction}
A finite dimensional algebra $\Lambda$ is called representation finite if there are only finitely many non-isomorphic finitely generated indecomposable $\Lambda$-modules. Gabriel \cite{gabriel1973indecomposable} showed that the path algebra of a connected quiver is representation finite if and only if its underlying diagram is a Dynkin diagram of type ADE. Later this was generalised to species by Dlab and Ringel \cite{dlab1975algebras}, who showed that the tensor algebra of a connected species is representation finite if and only if its underlying diagram is a Dynkin diagram of ABCDEFG. For studying the module category in these cases we use the preprojective algebra introduced by \cite{Gelfand1979ModelAA}, and later by Dlab and Ringel \cite{Dlab_1980} in the species case. It is well known that the preprojective algebra $\Pi(S)$ of a species $S$ does not depend on the orientation $S$. Each orientation of $S$ corresponds to a natural grading on $\Pi(S)$, such that the degree zero part in this grading is isomorphic to the tensor algebra of $S$.

In the paper \cite{APR1979} the notion of APR tilting was introduced. Performing APR tilting on a representation finite algebra yields another representation finite algebra. Let $S$ be a representation finite species over a quiver $Q$ with diagram $\Delta$. Then performing APR tilting on a representation finite species at a sink in $Q$ turns into a source and vice versa. Since the underlying quiver is a tree, we can obtain every possible orientation in $\Delta$. The description provided by \cite{Dlab_1980} says that every orientation of $Q$ yields a grading for the preprojective algebra. Thus, for species, APR tilting corresponds to a change of grading on the preprojective algebra. This means that if $S$ and $S'$ are both representation finite species such that $\Pi(S) = \Pi(S')$, then $T(S')$ can be retrieved by applying successive APR tilting on $T(S)$.

We now turn our focus towards the higher analogue of representation finite algebras, introduced by Iyama and Oppermann in \cite{IO11nRFalgandnAPRtilt}, called $d$-representation finite algebras. An algebra $\Lambda$ with $\gldim \Lambda \le d$ is $d$-representation finite if there exists a $d$-cluster tilting $\Lambda$-module. Iyama and Oppermann introduced the operation $d$-APR tilting, which is a generalisation of the ordinary APR tilting, and proved that it preserves $d$-representation finiteness. In particular, for $d=2$ this means that $2$-APR tilting preserves $2$-representation finiteness. In contrast to the case $d=1$ there does not exist any known classification of $d$-representation finite algebra for $d\ge 2$, not even for $d=2$. In \cite{IO13Stablecateofhigherpreproj} Iyama and Oppermann introduced the $(d+1)$-preprojective algebra $\Pi_{d+1}(\Lambda)$ as the higher analogue of the preprojective algebra whenever $\gldim \Lambda \le d$. Iyama and Oppermann proved that $\Lambda$ is $d$-representation finite if and only if $\Pi_{d+1}(\Lambda)$ is self-injective. In this paper we provide a proof that $d$-APR tilting on $\Lambda$ corresponds to a change of grading of $\Pi_{d+1}(\Lambda)$.

A partial classification of $2$-representation finite algebras over an algebraically closed field is due to Herschend and Iyama \cite{HerschendOsamu2011quiverwithpotential} using quivers with potentials. A quiver with potential is a pair $(Q, W)$ of a quiver $Q$ and a potential $W$, i.e. a sum of cycles in the path algebra of $Q$. Given a quiver with potential $(Q, W)$ we define the Jacobian algebra $\mathcal{P}(Q, W)$ as the path algebra of $Q$ modulo the ideal generated by all derivations of $W$. A quiver with potential is said to be self-injective if its corresponding Jacobian algebra is a finite dimensional self-injective algebra. A set of arrows $C$ in $Q$ is called a cut if it induces a grading on the arrows such that $W$ is homogeneous of degree $1$. For a cut $C$ the truncated Jacobian algebra $\mathcal{P}(Q, W, C)$ is defined as the path algebra over $Q_C$ modulo the ideal generated by all derivations of $W$ by elements from $C$, where $Q_C$ is the quiver obtained from $Q$ by removing all arrows in $C$. In \cite{HerschendOsamu2011quiverwithpotential} it was proved that every basic $2$-representation finite algebra over an algebraically closed field appears as a truncated Jacobian algebra of a self-injective quiver with potential. In this case we have that $\Pi_3(\mathcal{P}(Q, W, C)) \cong \mathcal{P}(Q, W)$. Herschend and Iyama also showed that performing $2$-APR tilting on $\mathcal{P}(Q, W, C)$ yields $\mathcal{P}(Q, W, C')$, where $C'$ is obtained from $C$ by performing an operation called cut-mutation, yielding a rich source of $2$-representation finite algebras. Furthermore, they also gave sufficient conditions for when all cuts are transitive under successive cut-mutation. Moreover, cut-mutation corresponds to a change of grading on $\mathcal{P}(Q, W)$.

In this paper we aim to generalise this correspondence between $2$-APR tilting and cut-mutation. We also provide sufficient conditions when all cuts are transitive under successive cut-mutation. For this we study a more general version of Jacobian algebras $\mathcal{P}(S, W)$ constructed using species with potential $(S, W)$, studied by \cite{Nquefack2012PotentialsJacobian}. We assume throughout the article that $\K$ is a perfect field. Let $(S, W)$ be a species with potential where $S$ is a species over a quiver $Q$ and $W$ a potential. Similarly, to the quiver case, we define cuts $C$ and truncated Jacobian algebras $\mathcal{P}(S, W, C)$. Our main result is a generalisation of \cite[Theorem 8.7]{HerschendOsamu2011quiverwithpotential}.

\begin{mythm*}(Theorem~\ref{theorem - jacobian 2apr tils of each other})
	Let $(S, W)$ be a self-injective simply connected species with potential with enough cuts. Assume that $(S, W)$ has a preprojective cut. Then all cuts $C$ are preprojective and the corresponding truncated Jacobian algebras $\mathcal{P}(S, W, C)$ are iterated $2$-APR tilts of each other. In particular, each $\mathcal{P}(S, W, C)$ is $2$-representation finite.
\end{mythm*}

Here a cut $C$ is said to be preprojective if $\gldim \mathcal{P}(S, W, C)\le 2$ and there is an isomorphism of complete graded algebras $\varphi: \Pi_3(\mathcal{P}(S, W, C)) \cong \mathcal{P}(S, W)$ with $\mathcal{P}(S, W)$ graded by $g_C$ such that $\varphi$ is the identity on the degree $0$ part. In the quiver case, all cuts for self-injective quivers with potentials are preprojective due to \cite{Keller2011CYcompletions}, see also \cite[Proposition 3.9]{HerschendOsamu2011quiverwithpotential}. By earlier work \cite{soderberg2022preprojective} we have a complete list of species for $S^1$ and $S^2$ such that $\Lambda=T(S^1)\otimes_\K T(S^2)$, i.e. the tensor product of tensor algebras of species, becomes a $2$-representation finite $\K$-algebra. In \cite{soderberg2024mutation} we showed that $\Lambda$ is not necessarily basic, but it is Morita equivalent to a basic $\K$-algebra $\Lambda'$. Thus we replace $\Lambda$ by $\Lambda'$ whenever $\Lambda$ is not a basic $\K$-algebra, in order to perform $2$-APR tilting. We prove in this case that there is a self-injective simply connected species with potential $(S, W)$ that has enough cuts such that $\Pi_3(\Lambda) = \mathcal{P}(S, W)$. We show that there exists a cut $C$ such that $\mathcal{P}(S, W, C)\cong \Lambda$, so in particular $C$ is a preprojective cut. Thus the above theorem applies and every truncated Jacobian algebra corresponding to a cut is a $2$-representation finite algebra.

The outline of the proof of our main theorem goes as follows. We introduce the notion of quivers with cycles, which is defined to be a quiver $Q$ together with a set of cycles $Q_2$. Every species with potential $(S, W)$ gives rise to a quiver with cycles, where $Q_2$ is given as the support of $W$. In particular, cuts can be defined on the level of quiver with cycles. We define an operation called cut-mutation following \cite{HerschendOsamu2011quiverwithpotential}. We give sufficient conditions on quivers with cycles for successive cut-mutations to act transitively on all cuts. Then the main theorem follows by showing that cut-mutation corresponds to $2$-APR tilting. Note that for any cut $C'$ reached by $2$-APR tilting we have that $\Pi_3(\mathcal{P}(S, W, C')) \cong \mathcal{P}(S, W)$ and $\mathcal{P}(S, W, C')$ is $2$-representation finite. Now the transitivity of cut-mutation implies that all truncated Jacobian algebras are $2$-representation finite.

Let us also demonstrate an example. We will refer to Example~\ref{ex - example for intro} for more details. We obtain a species with potential $(S, W)$, which is computed by taking the tensor product of two species of type $A_3$ and $B_2$ respectively, and is described by the following diagram.
\begin{equation*}
	\begin{tikzcd}[row sep = 0.6cm, column sep = 1cm]
		\C \arrow[r, "\C"] \arrow[dd, "\C"] & \C \arrow[dd, "\C"] & \C \arrow[l, "\C"'] \arrow[dd, "\C"] \\ \\
		\R \arrow[r, "\R"] & \R \arrow[luu, "\C"'] \arrow[ruu, "\C"'] & \R \arrow[l, "\R"']
	\end{tikzcd}
\end{equation*}
The support of the potential $W$ consists of all $3$-cycles. By earlier work \cite{soderberg2022preprojective} and \cite{soderberg2024mutation}, the Jacobian algebra $\mathcal{P}(S, W)$ is equal to $\Pi_3(\mathcal{P}(S, W, C))$ where $C$ is the cut consisting of the two diagonal arrows. Applying our main theorem now implies that any other cut yields a $2$-representation finite algebra. For example, we have the following $2$-representation finite algebra
\begin{equation*}
	\begin{tikzcd}[row sep = 0.6cm, column sep = 1cm]
		\C \arrow[r, "\C"] \arrow[dd, "\C"] & \C \arrow[dd, dotted] & \C \arrow[l, "\C"'] \arrow[dd, dotted] \\ \\
		\R \arrow[r, dotted] & \R \arrow[luu, "\C"'] \arrow[ruu, "\C"'] & \R \arrow[l, "\R"']
	\end{tikzcd}
\end{equation*}
where the chosen cut is represented by the dotted arrows and the relations given by derivation of the potential by the elements at the cut. \\

\noindent\textbf{Acknowledgements.} The author wishes to thank his supervisor Martin Herschend for all his guidance and helpful discussions while writing this article.

\subsection{Outline}
In Section~\ref{Section - prel} we establish preliminaries on species with potential and higher preprojective algebras. In Section~\ref{section - quiver with cycles} we define quivers with cycles, which are the underlying structures of a species with potentials. We define cuts for quivers with cycles and cut-mutation. We provide sufficient conditions for a quiver with cycles such that all cuts are transitive under successive cut-mutation. In Section~\ref{section - tensor products of quiver with cycles} we define quivers with cycles related to preprojective algebras of tensor products. We provide a set of example species with potential described using tensor products where all cuts are transitive under successive cut-mutation. In Section~\ref{section - truncated jacobian algebras} we define truncated Jacobian algebras. We introduce the notion of algebraic cuts for species with potentials. We prove that all cuts for every self-injective species with potential are algebraic. In Section~\ref{Section - Preproj, 2apr and cutmutation} we show that cut-mutation corresponds to $2$-APR tilting and present the main theorem of this article. We compute various examples of $2$-representation finite $\K$-algebras by applying our main theorem.

\section{Preliminaries}\label{Section - prel}
We will use the same setup and conventions as in \cite{soderberg2024mutation}. For convenience of the reader we will state the most important steps and will refer to \cite{soderberg2024mutation} for the details.

\subsection{Conventions}
In this article $R$-modules mean left $R$-modules and right $R$-modules are considered as left $R^{op}$-modules. All graded rings $R$ are assumed to be non-negatively graded, i.e. $R = \bigoplus_{i\ge 0}R_i$, where $R_iR_j \subseteq R_{i+j}$. We say that $R$ is complete graded if instead $R = \prod_{i\ge 0}R_i$, i.e. $R$ is the completion of the subring $\bigoplus_{i\ge 0}R_i$. The composition of arrows and morphisms goes from right to left. The center of an algebra $D$ is denoted by $\mathcal{Z}(D)$. The $D$-center $\mathcal{Z}_D(A)$, of some $D$-$D$-bimodule $A$, is defined to be the $\mathcal{Z}(D)$-subbimodule of $A$
\begin{equation*}
	\mathcal{Z}_D(A) = \{a\in A \mid ad = da, d\in D \}.
\end{equation*}
All quivers are assumed to be finite and connected.

\subsection{Casimir Elements}
The following lemma is well-known and we will refer to \cite[Lemma 4.1]{kulshammer2017pro} for the proof. It can also be found in \cite{PareigisLnotes}.

\begin{mylemma}\label{lemma - dual basis}(Dual basis lemma)
	Let $P_R$ be a right $R$-module. Then, the following statements are equivalent:
	\begin{enumerate}
		\item $P$ is finitely generated and projective.
		\item There are $x_1, x_2, \dots, x_m\in P$ and $f_1, f_2, \dots, f_m\in \mathrm{Hom}_{R^{op}}(P, R)$ such that for every $x\in P$ we have
		\begin{equation*}
			x = \sum_{i=1}^m x_if_i(x).
		\end{equation*}
		\item The dual basis homomorphism $P\otimes_R \mathrm{Hom}_{R^{op}}(P, R) \to \mathrm{End}_{R^{op}}(P)$, $p\otimes f\mapsto (q\mapsto pf(q))$ is an isomorphism.
	\end{enumerate}
\end{mylemma}

\begin{mydef}
	Let $x_i$ and $f_i$ be as in Lemma~\ref{lemma - dual basis}. The element
	\begin{equation*}
		\sum_{i=1}^n x_i\otimes_R f_i\in P\otimes_R \mathrm{Hom}_{R^{op}}(P, R)
	\end{equation*}
	is called the Casimir element of $P\otimes_R \mathrm{Hom}_{R^{op}}(P, R)$. The Casimir element does not depend on the choice of $x_i$ and $f_i$ by \cite[Lemma 1.1]{Dlab_1980}.
\end{mydef}

\subsection{Species with Potential}
Let $\K$ be a field.

\begin{mydef}\label{Definition - Species}(Species)
	Let $Q$ be a quiver without multiple arrows.
	\begin{enumerate}
		\item A species $S=(D_i, M_\alpha)_{i\in Q_0, \alpha\in Q_1}$ is a collection where each $D_i$ is a semisimple $\K$-algebra and $M_\alpha\in D_j$-$D_i$-$\mathrm{mod}$, where $\alpha:i\rightarrow j$, such that $\mathrm{Hom}_{D_i^{op}}(M_\alpha, D_i)\cong \mathrm{Hom}_{D_j}(M_\alpha, D_j)$ as $D_i$-$D_j$-modules, such that $\mathcal{Z}_\K(D_i) = D_i$ and $\mathcal{Z}_\K(M_\alpha) = M_\alpha$ for all $i\in Q_0$ and $\alpha\in Q_1$.
		\item A finite dimensional species $S$ is a species $S$ such that $\dim_\K D_i<\infty$ and $\dim_\K M_\alpha<\infty$ for all $i\in Q_0$ and $\alpha\in Q_1$.
		\item We say that a finite dimensional species $S$ is a $\mathbb{K}$-species if all $D_i$ are division $\K$-algebras over a common central subfield $\mathbb{K}$.
		\item We call a species $S$ acyclic if $Q$ is acyclic.
	\end{enumerate}
\end{mydef}

\begin{mydef}
	For a species $S$ let $D=\bigoplus_{i\in Q_0}D_i$ and $M=\bigoplus_{\alpha\in Q_1}M_\alpha\in D$-$D\mathrm{-mod}$. 
	\begin{enumerate}
		\item We define the tensor algebra $T(S)$ to be the tensor ring $T(D, M)$. More explicitly,
		\begin{equation*}
			T(S) = T(D, M) = \bigoplus_{k\ge 0} M^k = \bigoplus_{k\ge 0} M^{\otimes_D k}, \quad M^0 = D.
		\end{equation*}
		\item We define the complete tensor algebra $\widehat{T}(S)$ as
		\begin{equation*}
			\widehat{T}(S) = \prod_{k\ge 0} M^k = \prod_{k\ge 0} M^{\otimes_D k}, \quad M^0 = D.
		\end{equation*}
	\end{enumerate}
\end{mydef}

For a species $S$ the semisimple $\K$-algebra $D$ is a symmetric algebra due to \cite[Corollary 5.17]{skowronski2011frobenius}, i.e. there exists a morphism $\t: D\to \K$ such that $\t(ab) = \t(ba)$ for all $a,b\in D$ and that there are no non-zero ideals contained in the kernel of $\t$. As noted in \cite[Remark 2.2]{Nquefack2012PotentialsJacobian} the dualising condition, i.e. $\mathrm{Hom}_{D_i^{op}}(M_\alpha, D_i)\cong \mathrm{Hom}_{D_j}(M_\alpha, D_j)$ as $D_i$-$D_j$-modules, in Definition~\ref{Definition - Species} is automatic as $D$ is symmetric. Moreover, the $D$-$D$-bimodule morphisms
\begin{equation*}
	\begin{aligned}
		\t_l &= \t\circ -: \mathrm{Hom}_D(M, D) \xrightarrow{\sim} \mathrm{Hom}_\K(M, \K) \\
		\t_r &= \t\circ -: \mathrm{Hom}_{D^{op}}(M, D) \xrightarrow{\sim} \mathrm{Hom}_\K(M, \K)
	\end{aligned}
\end{equation*}
are isomorphisms.

Let us denote $M^* = \mathrm{Hom}_\K(M, \K)$. Following \cite{Nquefack2012PotentialsJacobian} we define the morphism
\begin{equation*}
	\begin{aligned}
		\mathfrak{b}: M\otimes_D M^* \oplus M^*\otimes_D M &\to D, \\
		x\otimes x^* + y^*\otimes y &\mapsto \t_l^{-1}(x^*)(x) + \t_r^{-1}(y^*)(y).
	\end{aligned}
\end{equation*}

We will denote the morphism $\b$ as $\b_S$ to make it clear that it is defined for a given species $S$.

\begin{mylemma}\cite[Lemma 3.9]{soderberg2024mutation}
	The morphism $\mathfrak{b}$ satisfies:
	\begin{enumerate}
		\item $\t$ is a symmetrising trace for $\b$, i.e. $\t(\mathfrak{b}(x\otimes x^*)) = \t(\b(x^*\otimes x))$ for all $x\in M$ and $x^*\in M^*$.
		\item The morphisms
		\begin{equation*}
			\begin{aligned}
				\b_{1, l}:& M^* \to \mathrm{Hom}_D(M, D), \\
				&\hspace*{0.2cm}x^*\mapsto \b(-\otimes x^*), \\
				\b_{1, r}:& M \to \mathrm{Hom}_{D^{op}}(M^*, D), \\
				&\hspace*{0.2cm}x \mapsto \b(x\otimes -),
			\end{aligned}
		\end{equation*}
		are bijective. Or equivalently, the morphisms
		\begin{equation*}
			\begin{aligned}
				\b_{2, r}:& M^* \to \mathrm{Hom}_{D^{op}}(M, D), \\
				&\hspace*{0.2cm}x^*\mapsto \b(x^*\otimes -), \\
				\b_{2, l}:& M \to \mathrm{Hom}_D(M^*, D), \\
				&\hspace*{0.2cm}x \mapsto \b(-\otimes x),
			\end{aligned}
		\end{equation*}
		are bijective.
	\end{enumerate}
\end{mylemma}

Let
\begin{equation*}
	\sum_{i=1}^n \hat{x}_i\otimes x_i\in \mathrm{Hom}_D(M, D)\otimes_D M, \qquad \sum_{j=1}^m y_j\otimes \hat{y}_j\in M\otimes_D \mathrm{Hom}_{D^{op}}(M, D),
\end{equation*}
be Casimir elements. Setting $x^*_i = \b_{1, l}^{-1}(\hat{x}_i)$ and $y^*_j = \b_{2, r}^{-1}(\hat{y}_j)$ we get elements
\begin{equation}\label{eq - Casimir elements of M M^*}
	\sum_{i=1}^n x_i^*\otimes x_i\in M^*\otimes_D M, \qquad \sum_{j=1}^m y_j\otimes y_j^*\in M\otimes_D M^*,
\end{equation}
satisfying
\begin{equation}\label{eq - casimir elements equations ass. to b}
	\sum_{i=1}^n \b(x\otimes x_i^*)x_i = x = \sum_{j=1}^my_j\b(y_j^*\otimes x), \qquad \sum_{i=1}^n x_i^*\b(x_i\otimes \xi) = \xi = \sum_{j=1}^m\b(\xi\otimes y_j)y_j^*,
\end{equation}
for all $x\in M$ and $\xi\in M^*$.

\subsection{Jacobian Algebras}
In this subsection we define the Jacobian algebra for a species as in \cite{Nquefack2012PotentialsJacobian}.

\begin{mydef}\cite[Definition 3.5]{Nquefack2012PotentialsJacobian}
	Let $S$ be a species. We call elements in $\mathcal{Z}_D(\widehat{T}(S)_{\ge 2})$ potentials.
\end{mydef}

Throughout this article we assume that all potentials are finite, i.e. we assume that all potentials $W\in \mathcal{Z}_D(\widehat{T}(S)_{\ge 2})$ are such that $W\in T(S)$.

\begin{mydef}\cite[Proposition 3.2]{Nquefack2012PotentialsJacobian}\label{definition - derivative operators}
	Let $S$ be a species. We define the derivative operators as
	\begin{equation*}
		\begin{aligned}
			\partial^l: M^*\otimes_D M\otimes_D \widehat{T}(S) &\to \widehat{T}(S), \\
			\xi\otimes a\otimes b &\mapsto \partial^l_\xi(a\otimes b) = \b(\xi\otimes a)b, \\
			\partial^r: \widehat{T}(S) \otimes_D M\otimes_D M^* &\to \widehat{T}(S), \\
			a\otimes b\otimes \xi &\mapsto \partial^r_\xi(a\otimes b) = a\b(b\otimes \xi).
		\end{aligned}
	\end{equation*}
\end{mydef}

\begin{mylemma}\cite[Lemma 2.5, Proposition 3.2]{Nquefack2012PotentialsJacobian}
	Let $S$ be a species and let $W\in \mathcal{Z}_D(\widehat{T}(S)_{\ge 2})$ be a potential. We define the permutation operators $\varepsilon_l, \varepsilon_r: \mathcal{Z}_D(\widehat{T}(S)_{\ge 2})\to \mathcal{Z}_D(\widehat{T}(S)_{\ge 2})$ by
	\begin{equation*}
		\begin{aligned}
			\varepsilon_l(W) = \sum_{i=1}^m \partial^l_{x^*_i}(W)x_i, \\
			\varepsilon_r(W) = \sum_{i=1}^n y\partial^r_{y^*_i}(W),
		\end{aligned}
	\end{equation*}
	where $x_i, x^*_i, y_i, y^*_i$ make up the Casimir elements in \eqref{eq - Casimir elements of M M^*}. We define the cyclic permutation operator $\varepsilon_c: \mathcal{Z}_D(\widehat{T}(S)_{\ge 2})\to \mathcal{Z}_D(\widehat{T}(S)_{\ge 2})$ on the homogeneous element of degree $d+1$ by $\varepsilon_c(W) = \sum_{k=0}^d \varepsilon^k_l(W) = \sum_{k=0}^d \varepsilon^k_r(W)$.
\end{mylemma}

For each $\xi\in M^*$ we use the notation $\partial^l_\xi = \partial^l(\xi\otimes-)$ and $\partial^r_\xi = \partial(-\otimes \xi)$.

\iffalse
\begin{myrem}
	The derivative operators in Definition~\ref{definition - derivative operators} satisfy the following:
	\begin{equation*}
		\begin{aligned}
			\partial^l_\xi(d) &= 0 = \partial^r_\xi(d), \\
			\partial^l_\xi(mx) &= \partial^l_\xi(m)x, \\
			\partial^r_\xi(xm) &= x\partial^r_\xi(m),
		\end{aligned}
	\end{equation*}
	for all $d\in D$, $m\in M$ and $x\in \widehat{T}(S)$.
\end{myrem}

\begin{mylemma}\cite[Proposition 3.2]{Nquefack2012PotentialsJacobian}
	Let $S$ be a species and $W\in \mathcal{Z}_D(\widehat{T}(S)_{\ge 2})$.
	\begin{enumerate}
		\item The derivative operators in Definition~\ref{definition - derivative operators} satisfy $\partial^r_\xi(\varepsilon_l (W)) = \partial^l_\xi(W)$ and $\partial^l_\xi(\varepsilon_r(W)) = \partial^r_\xi(W)$ for $\xi\in M^*$.
		\item $\varepsilon_l\circ \varepsilon_r(W) = W = \varepsilon_r\circ \varepsilon_l(W)$. 
	\end{enumerate}
\end{mylemma}
\fi

\begin{mydef}
	We define the cyclic derivative operator
	\begin{equation*}
		\begin{aligned}
			\partial: M^*\otimes_\K \mathcal{Z}_D(\widehat{T}(S)_{\ge 2}) &\to \widehat{T}(S)_{\ge 1} \\
			\xi\otimes W &\mapsto \partial^l(\xi\otimes \varepsilon_c(W)) = \partial^r(\varepsilon_c(W)\otimes \xi).
		\end{aligned}
	\end{equation*}
\end{mydef}

\begin{mydef}
	For a species with potential $(S, W)$ we define the Jacobian algebra $\mathcal{P}(S, W) = \widehat{T}(S)/\mathcal{J}(S, W)$, where $\mathcal{J}(S, W) = \overline{\langle \partial_\xi (W) \mid \xi\in M^*\rangle}$. 
\end{mydef}

\subsection{\textit{d}-Representation Finite Algebras}
The main part of this paper is to understand the relation between $2$-APR tilting and cut mutation defined in later sections. For this we will work with $2$-representation finite algebras.

\begin{mydef}
	Let $\Lambda$ be a finite dimensional $\K$-algebra with $\mathrm{gldim}\Lambda \le d$ for some positive integer $d$. We say that $\Lambda$ is a $d$-representation finite if there exists a $d$-cluster tilting object $M\in \mathrm{mod} (\Lambda)$, i.e.
	\begin{equation*}
		\begin{aligned}
			\mathrm{add}(M) &= \{X\in \Lambda\mathrm{-mod} \mid \mathrm{Ext}^i_\Lambda(M, X) = 0 \mbox{ for any } 0 < i < d \} = \\
			&= \{X\in \Lambda\mathrm{-mod} \mid \mathrm{Ext}^i_\Lambda(X, M) = 0 \mbox{ for any } 0 < i < d \}.
		\end{aligned}
	\end{equation*}
\end{mydef}

For a $d$-representation finite algebra $\Lambda$ we define the $d$-Auslander-Reiten translations as
\begin{equation*}
	\begin{aligned}
		\tau_d &= \mathrm{Tor}^\Lambda_d(\DD\Lambda, -)\cong \DD\mathrm{Ext}^d_\Lambda(-, \Lambda): \Lambda\mathrm{-mod}\rightarrow \Lambda\mathrm{-mod}, \\
		\tau_d^{-1} &= \DD\mathrm{Tor}^\Lambda_d(\DD-, \DD\Lambda)\cong \mathrm{Ext}^d_\Lambda(\DD\Lambda, -): \Lambda\mathrm{-mod}\rightarrow \Lambda\mathrm{-mod},
	\end{aligned}
\end{equation*}
where $\DD = \Hom_\K(-, \K)$.

\begin{mydef}\cite[Definition 2.11]{IO13Stablecateofhigherpreproj}
	Let $\Lambda$ be a finite dimensional algebra with $\gldim \Lambda \le d$. We denote the $(d+1)$-preprojective algebra, or the higher preprojective algebra, of $\Lambda$ by $\Pi_{d+1}(\Lambda)$. It is defined as
	\begin{equation*}
		\Pi_{d+1}(\Lambda) = \bigoplus_{i\ge 0}(\Pi_{d+1}(\Lambda))_i,
	\end{equation*}
	where
	\begin{equation*}
		\Pi_{d+1}(\Lambda)_i = \Pi_i = \mathrm{Hom}_\Lambda(\Lambda, \tau_d^{-i}\Lambda)
	\end{equation*}
	with multiplication
	\begin{equation*}
		\begin{aligned}
			\Pi_i\times \Pi_j&\to \Pi_{i+j}, \\
			(u, v)&\mapsto uv = (\tau^{-i}_d(v)\circ u: \Lambda\to \tau_d^{-(i+j)}\Lambda).
		\end{aligned}
	\end{equation*}
	The complete $(d+1)$-preprojective algebra of $\Lambda$ is defined as the completion of $\Pi_{d+1}(\Lambda)$, i.e.,
	\begin{equation*}
		\widehat{\Pi}_{d+1}(\Lambda) = \prod_{i\ge 0} \Hom_\Lambda(\Lambda, \tau^{-i}_d\Lambda).
	\end{equation*}
\end{mydef}

\begin{mydef}
	Let $\Lambda$ be a finite dimensional $\K$-algebra with $\gldim \Lambda < \infty$. We define $\nu_d = \nu\circ [-d]: \mathcal{D}^b(\Lambda\mathrm{-mod})\to \mathcal{D}^b(\Lambda\mathrm{-mod})$, where $\nu = D\Lambda\otimes_\Lambda^\mathbb{L}-$ the is the Nakayama functor.
\end{mydef}

\begin{myprop}
	For a finite dimensional algebra $\Lambda$ with $\gldim \Lambda \le d$, we can also define the $(d+1)$-preprojective algebra using $\nu_d^{-1}$ instead of $\tau_d^-$ since
	\begin{equation}\label{eq - derived prop preprojective}
		\Hom_{\mathcal{D}^b(\Lambda\mathrm{-mod})}(\Lambda, \tau_d^{-i}\Lambda) = \tau_d^{-i}\Lambda = H^0(\nu_d^{-i}\Lambda) = \Hom_{\mathcal{D}^b(\Lambda\mathrm{-mod})}(\Lambda, \nu_d^{-i}\Lambda),
	\end{equation}
	i.e.
	\begin{equation}\label{Remark - definition of derived preprojective algebra}
		\Pi_{d+1}(\Lambda) = \bigoplus_{i\ge 0}\mathrm{Hom}_{\mathcal{D}^b(\Lambda\mathrm{-mod})}(\Lambda, \nu_d^{-i}\Lambda),
	\end{equation}
	and similarly the complete $(d+1)$-preprojective algebra can be defined as
	\begin{equation*}
		\widehat{\Pi}_{d+1}(\Lambda) = \prod_{i=0}^\infty \Hom_{\mathcal{D}^b(\Lambda\mathrm{-mod})}(\Lambda, \nu^{-i}_d\Lambda).
	\end{equation*}
\end{myprop}

\begin{proof}
	The first and the third equality in \eqref{eq - derived prop preprojective} follows since
	\begin{equation*}
		\Hom_{\mathcal{D}^b(\Lambda\mathrm{-mod})}(\Lambda, X) = H^0(X)
	\end{equation*}
	for $X\in \mathcal{D}^b(\Lambda\mathrm{-mod})$ and the second equality in \eqref{eq - derived prop preprojective} follows from \cite[Lemma 5.5 (b)]{Iyama2011clustertiltinghigheraus}.
\end{proof}

Let $\Lambda$ be a basic $d$-representation finite $\K$-algebra. Let $\{e_i\}_{i\in I}$ be a complete set of primitive orthogonal idempotents. Then by \cite[Proposition 1.3 (b)]{Iyama2011clustertiltinghigheraus} there exists a permutation $\sigma:I\to I$, which is called the Nakayama permutation, such that
\begin{equation*}
	\Lambda e_{\sigma(i)} = \tau_d^{l_i - 1}\DD(e_i\Lambda)
\end{equation*}
for some integers $l_i$.

\begin{mydef}\cite[Definition 1.2]{herschend2011n}
	A basic $d$-representation finite $\K$-algebra $\Lambda$ is said to be $l$-homogeneous if $l_i = l$ for all $i$.
\end{mydef}

We have the following useful result.

\begin{mycor}\cite[Corollary 1.5]{herschend2011n}\label{corollary - A x B is n_1 + n_2 RF}
	Let $\mathbb{K}$ be a perfect field and $l$ a positive integer. If $\Lambda_i$ is $l$-homogeneous $d_i$-representation finite for each $i\in \{1, 2\}$, then $\Lambda_1\otimes_\K \Lambda_2$ is an $l$-homogeneous $(d_1 + d_2)$-representation-finite algebra with a $(d_1 + d_2)$-cluster tilting module $\bigoplus_{i=0}^{l-1}(\tau_{d_1}^{-i}\Lambda_1\otimes_\K \tau_{d_2}^{-i}\Lambda_2)$.
\end{mycor}

\begin{myrem}\label{remark - nakayama permutation tensor product}
	In the situation of Corollary~\ref{corollary - A x B is n_1 + n_2 RF} we have that
	\begin{equation*}
		\tau_{d_1 + d_2}^{-i}(\Lambda e_s \otimes_\K \Lambda_2 e_t) = \tau_{d_1}^{-i}(\Lambda_1 e_s)\otimes_\K \tau_{d_2}^{-i}(\Lambda_2 e_t)
	\end{equation*}
	for $0\le i\le l$, $s\in I_1$ and $t\in I_2$. Taking $i=l$ gives the Nakayama permutation $\sigma: I_1\times I_2\to I_1\times I_2$, sending $\sigma(s, t) = (\sigma_1(s), \sigma_2(t))$, where $\sigma_1$ and $\sigma_2$ are the Nakayama permutations for $\Lambda_1$ and $\Lambda_2$ respectively.
\end{myrem}

The classification of (1-)representation finite species is done by \cite{dlab1975algebras}. It is stated as the following theorem.

\begin{mythm}\cite[Theorem B]{dlab1975algebras}\label{theorem - dlab ringel}
	A species $S$ is representation finite if the diagram of $S$ is a Dynkin diagram.
\end{mythm}

Let $\K\subseteq F\subseteq G$ be division $\K$-algebras with the valuation $k=\dim_\K G / \dim_\K F$. All representation finite species are given in Figure~\ref{Figure - Dynkin Diagrams} with their diagrams shown in Figure \ref{Figure - Dynkin Diagrams 1}. We omit writing out the orientation in Figure~\ref{Figure - Dynkin Diagrams} since every possible orientation is representation finite species with the same diagram for each type.

\begin{figure}[ht]
	\vspace{-0.5em}\begin{equation*}
		\begin{aligned}
			&\begin{tikzcd}
			(A_n) & F \arrow["F", -, r] & F \arrow["F", r, -] & \cdots \arrow["F", r, -] & F \arrow["F", r, -] & F \\
			\end{tikzcd} \\[-2em]
			&\begin{tikzcd}
			(B_n) & F & G \arrow["G"', -, l] \arrow["G", r, -] & \cdots \arrow["G", r, -] & G \arrow["G", r, -] & G \\
			\end{tikzcd} \\[-2em]
			&\begin{tikzcd}
			(C_n) & G \arrow["G", -, r] & F \arrow["F", r, -] & \cdots \arrow["F", r, -] & F \arrow["F", r, -] & F \\
			\end{tikzcd} \\[-2em]
			&\begin{tikzcd}[row sep = 0em]
			& F \\
			(D_n) & & F \arrow["F", -, r] \arrow["F"', lu, -] \arrow["F"', ld, -] & \cdots \arrow["F", r, -] & F \arrow["F", r, -]& F \\
			&F \\
			\end{tikzcd} \\[-0.5em]
			&\begin{tikzcd}[row sep=0em]
			& F \arrow["F", r, -] & F \arrow["F", r, -] & F \arrow["F", r, -] \arrow["F", dd, -] & F \arrow["F", r, -] & F \\
			(E_6) \\
			& & & F \\
			\end{tikzcd} \\[-0.3em]
			&\begin{tikzcd}[row sep=0em]
			& F \arrow["F", r, -] & F \arrow["F", r, -] & F \arrow["F", r, -] \arrow["F", dd, -] & F \arrow["F", r, -] & F \arrow["F", r, -] & F \\
			(E_7) \\
			& & & F \\
			\end{tikzcd} \\[-0.3em]
			&\begin{tikzcd}[row sep=0em]
			& F \arrow["F", r, -] & F \arrow["F", r, -] & F \arrow["F", r, -] \arrow["F", dd, -] & F \arrow["F", r, -] & F \arrow["F", r, -] & F \arrow["F", r, -] & F \\
			(E_8) \\
			& & & F \\
			\end{tikzcd} \\[-0.3em]
			&\begin{tikzcd}
			(F_4) & G \arrow["G", r, -] & G \arrow["G", r, -] & F \arrow["F", r, -] & F \\
			\end{tikzcd} \\[-2em]
			&\begin{tikzcd}
			(G_2) & G \arrow["G", r, -] & F
			\end{tikzcd}
		\end{aligned}
	\end{equation*}
	\caption{Description of species over Dynkin diagrams}\label{Figure - Dynkin Diagrams}
\end{figure}

\begin{figure}[ht]
	\vspace{-0.5em}\begin{equation*}
		\begin{aligned}
		&\begin{tikzcd}
		(A_n) & 1 \arrow[-, r] & 2 \arrow[r, -] & \cdots \arrow[r, -] & n-1 \arrow[r, -] & n \\
		\end{tikzcd} \\[-2em]
		&\begin{tikzcd}
		(B_n) & 1 & 2 \arrow["(2)"', l, Rightarrow] \arrow[r, -] & \cdots \arrow[r, -] & n-1 \arrow[r, -] & n \\
		\end{tikzcd} \\[-2em]
		&\begin{tikzcd}
		(C_n) & 1 \arrow["(2)", r, Rightarrow] & 2 \arrow[r, -] & \cdots \arrow[r, -] & n-1 \arrow[r, -] & n \\
		\end{tikzcd} \\[-2em]
		&\begin{tikzcd}[row sep = 0em]
		& 1 \\
		(D_n) & & 3 \arrow[-, r] \arrow[lu, -] \arrow[ld, -] & \cdots \arrow[r, -] & n-1 \arrow[r, -]& n \\
		&2 \\
		\end{tikzcd} \\[-0.5em]
		&\begin{tikzcd}[row sep=0em]
		& 1 \arrow[r, -] & 2 \arrow[r, -] & 3 \arrow[r, -] \arrow[dd, -] & 4 \arrow[r, -] & 5 \\
		(E_6) \\
		& & & 6 \\
		\end{tikzcd} \\[-0.3em]
		&\begin{tikzcd}[row sep=0em]
		& 1 \arrow[r, -] & 2 \arrow[r, -] & 3 \arrow[r, -] \arrow[dd, -] & 4 \arrow[r, -] & 5 \arrow[r, -] & 6 \\
		(E_7) \\
		& & & 7 \\
		\end{tikzcd} \\[-0.3em]
		&\begin{tikzcd}[row sep=0em]
		& 1 \arrow[r, -] & 2 \arrow[r, -] & 3 \arrow[r, -] \arrow[dd, -] & 4 \arrow[r, -] & 5 \arrow[r, -] & 6 \arrow[r, -] & 7 \\
		(E_8) \\
		& & & 8 \\
		\end{tikzcd} \\[-0.3em]
		&\begin{tikzcd}
		(F_4) & 1 \arrow[r, -] & 2 \arrow["(2)", r, Rightarrow] & 3 \arrow[r, -] & 4 \\
		\end{tikzcd} \\[-2em]
		&\begin{tikzcd}
		(G_2) & 1 \arrow["(3)", r, Rightarrow] & 2
		\end{tikzcd}
		\end{aligned}
	\end{equation*}
	\caption{Dynkin Diagrams}\label{Figure - Dynkin Diagrams 1}
\end{figure}

The relation between $l$-homogeneous species and representation finite species is described in the following corollary. 

\begin{mycor}\label{Corollary - l homogeneous species}\cite[Corollary 3.7]{soderberg2022preprojective}
	Let $S$ be a representation finite species. Then $S$ is $l$-homogeneous if $Q$ is stable under $\sigma$. Moreover, for the different cases, the integer $l$ is
	\begin{center}
		\begin{tabular}{c|c|c|c|c|c|c|c|c|c}
			$\Delta$ & $A_n$ & $B_n$ & $C_n$ & $D_n$ & $E_6$ & $E_7$ & $E_8$ & $F_4$ & $G_2$ \\\hline
			$l$ & $\frac{n+1}{2}$ & $n$ & $n$ & $n-1$ & $6$ & $9$ & $15$ & $6$ & $3$
		\end{tabular}
	\end{center}
\end{mycor}

\begin{myrem}
	The Nakayama permutation can be extended uniquely from the vertices to $Q$ by \cite[Lemma 3.3]{soderberg2022preprojective}. The Nakayama permutation is given by \cite[Theorem 3.1]{soderberg2022preprojective}, in particular it is equal to the identity morphism when the diagram of $Q$ is non-simply laced or is $D_{2n}$ for some integer $n$. Moreover, the Nakayama automorphism of the corresponding preprojective algebra is explicitly given in \cite[Theorem 5.3]{soderberg2022preprojective}.
\end{myrem}

The upshot of Corollary~\ref{Corollary - l homogeneous species} is that if we combine it with Corollary~\ref{corollary - A x B is n_1 + n_2 RF} we get many examples of $2$-representation finite $\K$-algebras. Furthermore, in these cases their $3$-preprojective algebra can be described up to Morita equivalence as Jacobian algebras of a species with potential, see Proposition~\ref{proposition - preprojective algebra of tensor species}. By \cite[Proposition 10.11]{soderberg2022preprojective} we have a description of the Nakayama permutation, and Nakayama automorphism, of said examples, see Remark~\ref{remark - nakayama permutation tensor product}.

\section{Quiver with Cycles}\label{section - quiver with cycles}
In this section we introduce the notion of quiver with cycles which is inspired by \cite[Section 8]{HerschendOsamu2011quiverwithpotential} and we will restate the definitions in \cite[Section 6 and Section 7]{HerschendOsamu2011quiverwithpotential} adjusted to the setting of quivers with cycles. In Section~\ref{Section - Preproj, 2apr and cutmutation} we show how cuts correspond to $2$-APR-tilts of the Jacobian algebra. In this section we study the combinatorics of cuts. We use the results of \cite{HerschendOsamu2011quiverwithpotential} in the setting when the proofs only depend on the quiver with cycles and not the underlying potential.

\subsection{Quiver with Cycles and Cuts}
A cycle in a quiver $Q$ is a path that starts and ends at the same vertex.

\begin{mydef}
	We call the triple $Q = (Q_0, Q_1, Q_2)$ quiver with cycles, where $(Q_0, Q_1)$ is a quiver and $Q_2$ is a set of cycles in $(Q_0, Q_1)$. 
\end{mydef}

The main motivation for studying quivers with cycles stems from the following definition.

\begin{mydef}
	Let $(S, W)$ be a species with potential, where $S$ is a species over $Q$. We define the quiver with cycles associated to $(S, W)$, also denoted by $Q$, as a quiver with cycles where we let $Q_2$ to be a minimal set of paths in $Q$ such that
	\begin{equation*}
		W \in \bigoplus_{\alpha_n\alpha_{n-1}\cdots \alpha_1\in Q_2}\bigoplus_{i = 1}^nM_{\alpha_i}\otimes_D M_{\alpha_{n-1}}\otimes_D \dots \otimes_D M_{\alpha_1}\otimes_D M_{\alpha_n} \otimes_D \dots \otimes_D M_{\alpha_{i+1}}.
	\end{equation*}
\end{mydef}

Whenever $(S, W)$ is a species with potential, $Q$ is assumed to be the quiver with cycles defined as above.

\begin{myrem}
	Note that $Q_2$ is unique up to choosing a starting vertex for each $c\in Q_2$. The results in this article using $Q_2$ do not depend on the choice of minimal paths. Therefore, we will no longer mention the choice of minimal set of paths.
\end{myrem}

For a quiver with cycles $Q$ and a subset $C\subseteq Q_1$ we call the pair $(Q, C)$ a truncated quiver with cycles. In this case $C$ introduces a grading on $Q_1$ via
\begin{equation*}
	g_C(\alpha) = \begin{cases}
		1, & \text{if }\alpha\in C \\
		0, & \text{if }\alpha\in Q_1\backslash C
	\end{cases}
\end{equation*}
which naturally extends to paths by taking $g_C(\alpha_1 \alpha_2 \cdots \alpha_n) = g_C(\alpha_1) + g_C(\alpha_2) + \dots + g_C(\alpha_n)$, where $\alpha_1\alpha_2\cdots \alpha_n$ is a path in $Q$. Moreover, if $g_C(c) = 1$ for all $c\in Q_2$ we say that $C$ is a cut of $Q$. We will adopt the notation that we represent elements in $C$ as dotted lines in $Q$.

Let $\overline{Q}$ be the doubled quiver of $Q$, i.e. $\overline{Q}_0 = Q_0$ and $\overline{Q}_1 = \{\alpha: i\to j, \alpha^{-1}:j\to i \mid \alpha\in Q_1 \}$. A walk in $Q$ is defined to be a path in $\overline{Q}$. We extend $g_C$ to walks by setting $g_C(\alpha^{-1}) = -g_C(\alpha)$ where $\alpha\in Q_1$. This means that
\begin{equation*}
	g_C(\alpha_1^{\epsilon_1}\alpha_2^{\epsilon_2}\cdots \alpha_n^{\epsilon_n}) = \sum_{i=0}^n \epsilon_i g_C(\alpha_i)
\end{equation*}
where $\epsilon_i \in \{-1, 1\}$ for all $0\le i\le n$.

\begin{myrem}
	Let $(S, W)$ be a species with potential. If $C\subseteq Q_1$ is a cut then $W$ is homogeneous of degree $1$ with respect to $g_C$.
\end{myrem}

\begin{mydef}\label{definition - cut-mutation}\cite[Definition 6.10]{HerschendOsamu2011quiverwithpotential}
	Let $Q$ be a quiver with cycles and $C\subseteq Q_1$.
	\begin{enumerate}[label=(a)]
		\item We say that a vertex $x\in Q_0$ is a strict source of $(Q, C)$ if all arrows ending at $x$ belong to $C$ and all arrows starting at $x$ do not belong to $C$.
		\item For a strict source $x\in Q_0$ of $(Q, C)$, we define the subset $\mu_x^+(C)\subseteq Q_1$ by removing all arrows in $Q$ ending at $x$ from $C$ and adding all arrows in $Q$ starting at $x$ to $C$.
		\item Dually we define a strict sink $x$ and $\mu_x^-(C)\subseteq Q_1$. In particular, $\mu_x^-\circ \mu_x^+(C) = C$ and $\mu_x^+\circ \mu_x^-(C) = C$ whenever is defined.
	\end{enumerate}
\end{mydef}

We call these operations cut-mutation.

\begin{myrem}
	Note that for a cut $C\subseteq Q_1$ and a strict source $x\in Q_1$ we have $g_C(p) = g_{\mu^+_x(C)}(p)$ for all cycles $p$ in $Q$. In particular, this means that $\mu^+_x(C)$ is also a cut. Similarly, if $x$ is a strict sink then $\mu^-_x(C)$ is a cut.
\end{myrem}

\begin{mydef}\cite[Definition 6.5, Definition 6.15, Definition 7.4]{HerschendOsamu2011quiverwithpotential}
	Let $Q$ be a quiver with cycles. We say that
	\begin{enumerate}
		\item two cuts $C, C'\subseteq Q_1$ are compatible if for any cyclic walk $p$ the equality $g_C(p) = g_{C'}(p)$ holds. In this case we write $C\sim C'$.
		\item $Q$ is fully compatible if all cuts are compatible with each other.
		\item $Q$ is covered if every arrow in $Q_1$ appears in a cycle in $Q_2$.
		\item $Q$ has enough cuts if each arrow in $Q$ is contained in a cut.
	\end{enumerate}
\end{mydef}

\begin{myprop}\label{proposition - Q full comp and enouch cuts}
	Let $Q$ be a quiver with cycles which is fully compatible and let $C$ be a cut. Then the following are equivalent:
	\begin{enumerate}
		\item $Q_C$ is an acyclic quiver.
		\item $Q$ has enough cuts.
	\end{enumerate}
\end{myprop}

\begin{proof}
	The statement that $Q_C$ is an acyclic quiver is equivalent to $Q_1 = \bigcup_{C'\sim C}C'$ by \cite[Proposition 6.16]{HerschendOsamu2011quiverwithpotential}. Now every cut satisfies $C'\sim C$ since $Q$ is fully compatible, so $Q_1 = \bigcup_{C'\sim C}C'$ means precisely that $Q$ has enough cuts.
\end{proof}

\begin{myrem}\label{remark - enouch cuts quiver acyclic}
	Note that if $Q$ has enough cuts, then by Proposition~\ref{proposition - Q full comp and enouch cuts} $Q_C$ is acyclic for every cut $C$.
\end{myrem}

\begin{mylemma}\label{lemma - Q suff cyclic and enouch comp}
	Let $Q$ be a quiver with cycles which is fully compatible, covered and has enough cuts and $C$ be a cut, then the following are true:
	\begin{enumerate}
		\item for each $\alpha\in Q_1$ there is a cycle $p$ containing $\alpha$ satisfying $g_C(p)= 1$.
		\item $Q_1 = \bigcup_{C \sim C'}C'$.
	\end{enumerate}
\end{mylemma}

\begin{proof}
	\begin{enumerate}
		\item Since $C$ is cut we have that $g_C(c) = 1$ for every $c\in Q_2$. Since $Q$ is covered, every $\alpha\in Q_1$ is contained in at least one cycle in $Q_2$.
		\item That $Q$ has enough cuts implies that every arrow is contained in a cut by definition and now the assertion follows by using that $Q$ is fully compatible.
	\end{enumerate}
\end{proof}

\begin{mythm}\label{theorem - qwc fully, covered enough implies trans}
	Let $Q$ be a quiver with cycles which is fully compatible, covered and has enough cuts, then the set of all cuts of $Q$ is transitive under successive cut-mutations.
\end{mythm}

\begin{proof}
	Lemma~\ref{lemma - Q suff cyclic and enouch comp} implies that $(Q, C)$ is sufficiently cyclic and has enough compatibles, in the sense of \cite[Definition 6.15]{HerschendOsamu2011quiverwithpotential}, for every cut $C$. The result then follows by applying \cite[Theorem 7.2]{HerschendOsamu2011quiverwithpotential}.
\end{proof}

\subsection{Canvas of a Quiver with Cycles}We will now introduce a CW complex for a quiver with cycles $Q$ as in \cite{HerschendOsamu2011quiverwithpotential} in order to give sufficient conditions for when $Q$ is fully compatible.

\begin{mydef}\cite[Definition 8.1]{HerschendOsamu2011quiverwithpotential}
	Let $Q$ be a quiver with cycles. We define the CW complex $X_Q$ as follows:
	\begin{enumerate}
		\item The $0$-cells are the vertices of $Q$, i.e. $X_Q^0 := Q_0$.
		\item The $1$-cells are indexed by the arrows $\alpha\in Q_1$ with attaching maps $\phi_\alpha: \{0, 1\}\to Q_0$ defined by $\phi_\alpha(0) = s(\alpha)$ and $\phi_\alpha(1) = t(\alpha)$.
		\item the $2$-cells are indexed by the cycles $c\in Q_2$ with attaching maps $\phi_c: S^1 \to X_S$ defined by
		\begin{equation*}
			\phi_c\left(\cos\left(\frac{2\pi}{s}(i+t) \right), \sin\left(\frac{2\pi}{s}(i+t) \right) \right)
		\end{equation*}
		for each integer $0\le i< s$ and a real number $0\le t<1$.
	\end{enumerate}
\end{mydef}

\begin{mydef}\cite[Definition 8.5]{HerschendOsamu2011quiverwithpotential}
	We say that a quiver with cycles $Q$ is simply connected if $X_Q$ is simply connected.
\end{mydef}

\begin{myrem}
	Note that if a quiver with cycles $Q$ is a tree, then $Q_2$ is empty. Moreover, it directly follows that $X_Q$ is simply connected.
\end{myrem}

\begin{myprop}\label{proposition - simply conn implies fully comp}\cite[Proposition 8.6]{HerschendOsamu2011quiverwithpotential}
	Every simply connected quiver with cycles is fully compatible.
\end{myprop}

\begin{myrem}
	Note that the proof of Proposition~\ref{proposition - simply conn implies fully comp} in \cite{HerschendOsamu2011quiverwithpotential} only depends on the quiver with cycle $Q$ and not the actual potential.
\end{myrem}

Combining Proposition~\ref{proposition - simply conn implies fully comp} and Theorem~\ref{theorem - qwc fully, covered enough implies trans} we get the following result.

\begin{mythm}\label{theorem - simply con implies trans}
	Let $Q$ be a simply connected quiver with cycles, which is covered and has enough cuts. Then the set of all cuts of $Q$ is transitive under successive cut-mutations.
\end{mythm}

Note that enough cuts can be checked by verifying that $Q_C$ is acyclic for any cut $C$ due to Proposition~\ref{proposition - Q full comp and enouch cuts}.

\section{Tensor Products and Quivers with Cycles}\label{section - tensor products of quiver with cycles}
In this section we provide a way of producing examples using tensor products where Theorem~\ref{theorem - simply con implies trans} applies. In \cite{soderberg2024mutation} we describe the $3$-preprojective algebra of the tensor product of two acyclic species using a species with potential, which is Morita equivalent to a $\K$-species with potential. In this case, under certain assumptions, we prove that all cuts are transitive under successive cut-mutations.

We start by constructing a quiver with cycles that corresponds to the $3$-preprojective algebra of the tensor product of two acyclic species.

\subsection{Quiver with Cycles for Tensor Product}
For a quiver $Q$ we define $Q^*$ to be the opposite quiver, i.e. $Q^*_0 = Q_0$ and $Q_1^* = \{\alpha^*: j\to i \mid \alpha:i\to j\in Q_1\}$. For quivers $Q^1$ and $Q^2$ we define $Q^1\tilde{\otimes} Q^2$ as the quiver with cycles given by
\begin{equation*}
	\begin{aligned}
		(Q^1\tilde{\otimes} Q^2)_0 &= Q^1_0\times Q^2_0, \\
  		(Q^1\tilde{\otimes} Q^2)_1 &= (Q^1_0\times Q^2_1) \coprod (Q^1_1\times Q^2_0)\coprod (Q^{1*}_1\times Q^{2*}_1), \\
  		(Q^1\tilde{\otimes}Q^2)_2 &= \{(\alpha^*, \beta^*)(\alpha, t(\beta))(s(\alpha), \beta), (\alpha^*, \beta^*)(t(\alpha), \beta)(\alpha, s(\beta)) \mid \alpha\in Q^1_1, \beta\in Q^2_1 \}.
	\end{aligned}
\end{equation*}
We extend the definition of source and target, $s$ and $t$, to be the identity on $Q_0$ for any given quiver $Q$, and define
\begin{equation*}
	s, t: (Q^1\tilde{\otimes} Q^2)_1 \to (Q^1\tilde{\otimes} Q^2)_0
\end{equation*}
by $s(\alpha_1, \alpha_2) = (s(\alpha_1), s(\alpha_2))$ and $t(\alpha_1, \alpha_2) = (t(\alpha_1), t(\alpha_2))$.

Note that the sets
\begin{equation}\label{eq - standard cuts}
	\begin{aligned}
		C_1 &= Q_0^1\times Q_1^2 \\
		C_2 &= Q_1^1 \times Q_0^2 \\
		C_3 &= (Q_1^1)^* \times (Q_1^2)^*
	\end{aligned}
\end{equation}
are three distinguished cuts for $Q^1\tilde{\otimes} Q^2$.

\begin{mylemma}\label{lemma - Q^1 x Q^2 covered enough cuts}
	The quiver with cycles $Q^1\tilde{\otimes}Q^2$
	\begin{enumerate}
		\item is covered,
		\item has enough cuts.
	\end{enumerate}
\end{mylemma}

\begin{proof}$\space$
	\begin{enumerate}
		\item Since every arrow in $(Q^1\tilde{\otimes}Q^2)_1$ is included in $(Q^1\tilde{\otimes}Q^2)_2$ by definition we have that $Q^1\tilde{\otimes}Q^2$ is covered.
		\item Every cycle in $(Q^1\tilde{\otimes}Q^2)_2$ consists of one arrow from $C_1$, one arrow from $C_2$ and one arrow from $C_3$. Since $C_1$, $C_2$ and $C_3$ are cuts it follows that $Q^1\tilde{\otimes}Q^2$ has enough cuts.
	\end{enumerate}
\end{proof}

To apply Theorem~\ref{theorem - simply con implies trans} we need to show that $Q^1\tilde{\otimes}Q^2$ is simply connected. This is not true in general and thus we need more assumptions.

\begin{mylemma}\label{lemma - finite tress product trans}
	For trees $Q^1$ and $Q^2$ the following are true.
	\begin{enumerate}
		\item $X_{Q^1\tilde{\otimes}Q^2}$ and $X_{Q^1}\times X_{Q^2}$ are homeomorphic.
		\item The set of all cuts on $Q^1\tilde{\otimes} Q^2$ is transitive under successive cut-mutations.
	\end{enumerate}
\end{mylemma}

\begin{proof}$\space$
	\begin{enumerate}
		\item We can view $Q^1$ and $Q^2$ as quivers with cycles by setting $Q^1_2 = Q^2_2 = \emptyset$. We will sketch the homeomorphism $\phi: X_{Q^1}\times X_{Q^2} \to X_{Q^1\tilde{\otimes}Q^2}$. By \cite[Chapter 0]{Hatcher2002algebraictopology} the CW complex of $X_{Q^1}\times X_{Q^2}$ is homeomorphic to the CW complex which
		\begin{enumerate}
			\item $0$-cells are the vertices $Q^1_0\times Q^2_0$,
			\item $1$-cells are indexed by the arrows $\{(\alpha, i) \mid i\in Q^1_0, \alpha\in Q^1_1\} \cup \{(j, \beta) \mid j\in Q^2_0, \beta\in Q^2_1\}$ attached at their endpoints,
			\item $2$-cells are indexed by $(\alpha, \beta)\in Q_1\times Q_2$ attached along the following diagrams.
			\begin{equation}\label{eq - 2 cell diagram}
				\begin{tikzcd}
					(s(\alpha){,} s(\beta)) \arrow[r, "(\alpha{,} s(\beta))"] \arrow[d, "(s(\alpha){,} \beta)"] & (t(\alpha){,} s(\beta)) \arrow[d, "(t(\alpha){,} \beta)"] \\
					(s(\alpha){,} t(\beta)) \arrow[r, "(\alpha{,} t(\beta))"] & (t(\alpha{,} t(\beta)))
				\end{tikzcd}
			\end{equation}
		\end{enumerate}
		with the natural attaching maps. Now let $\phi$ be the identity on the $0$-cells and $\phi$ be the inclusion on the $1$-cells. Each $2$-cell \eqref{eq - 2 cell diagram} is mapped by $\phi$ to the gluing of the two cells
		\begin{equation*}
			\begin{tikzcd}
				(s(\alpha){,} s(\beta)) \arrow[r, "(\alpha{,} s(\beta))"] \arrow[d, "(s(\alpha){,} \beta)"] & (t(\alpha){,} s(\beta)) \arrow[d, "(t(\alpha){,} \beta)"] \\
				(s(\alpha){,} t(\beta)) \arrow[r, "(\alpha{,} t(\beta))"] & (t(\alpha{,} t(\beta))) \arrow[lu, "(\alpha^*{,} \beta^*)", swap]
			\end{tikzcd}
		\end{equation*}
		in $X_{Q^1\tilde{\otimes}Q^2}$. It is straightforward to show that $\phi$ is indeed a continuous map with a continuous inverse.
		\item Since $Q^1$ and $Q^2$ are connected $X_{Q^1}$ and $X_{Q^2}$ are path connected. Thus applying \cite[Proposition 1.12]{Hatcher2002algebraictopology} we have that
		\begin{equation*}
			\pi_1(X_{Q^1}\times X_{Q^2}) \cong \pi_1(X_{Q^1})\times \pi_1(X_{Q^2}).
		\end{equation*}
		Now since $Q^1$ and $Q^2$ are trees $X_{Q^1}$ and $X_{Q^2}$ are simply connected. Hence (1) implies that $Q^1\tilde{\otimes}Q^2$ is simply connected.

		Together with Lemma~\ref{lemma - Q^1 x Q^2 covered enough cuts} we can apply Theorem~\ref{theorem - simply con implies trans} to get the result.
	\end{enumerate}
\end{proof}

\subsection{Tensor Products of Species}
In this section we recall results of \cite{soderberg2024mutation}. First we recall the description of the $3$-preprojective algebra of tensor products of species using species with potential. Then we recall that every species with potential is Morita equivalent to a $\K$-species with potential, i.e. their Jacobian algebras are Morita equivalent. Finally we show that, under certain conditions, the set of all cuts for said Morita equivalent $\K$-species with potential is transitive under successive cut-mutation.

Following \cite{soderberg2024mutation} we introduce the following notation. For every $\alpha\in Q_1$ we define $\underline{\alpha}\subseteq M_\alpha$ as sets of elements such that
\begin{equation*}
	c_\alpha = \sum_{a\in \underline{\alpha}} a \otimes a^* \in M_\alpha \otimes_{D_{s(\alpha)}} M_\alpha^*
\end{equation*}
is the Casimir element of $M_\alpha\otimes_{D_{s(\alpha)}}M_\alpha^*$ associated to $\b_S$, i.e.
\begin{equation*}
	\sum_{a\in \underline{\alpha}} a \otimes \b(a^*\otimes -)\in M_\alpha \otimes_{D_{s(\alpha)}} \mathrm{Hom}_{D_{s(\alpha)}^{op}}(M, D_{s(\alpha)})
\end{equation*}
is the Casimir element of $M_\alpha\otimes_{D_{s(\alpha)}} \mathrm{Hom}_{D_{s(\alpha)}^{op}}(M, D_{s(\alpha)})$. Similarly, we define $\overline{\alpha}\subseteq M_\alpha$ such that
\begin{equation*}
	c_{\alpha^*} = \sum_{a'\in \overline{\alpha}} {a'}^* \otimes a' \in M_\alpha^* \otimes_{D_{t(\alpha)}} M_\alpha
\end{equation*}
is the Casimir element of $M_\alpha^*\otimes_{D_{s(\alpha)}} M_\alpha$ associated to $\b_S$, i.e.
\begin{equation*}
	\sum_{a'\in \overline{\alpha}} \b(-\otimes {a'}^*)\otimes a'\in \mathrm{Hom}_{D_{t(\alpha)}}(M, D_{t(\alpha)}) \otimes_{D_{t(\alpha)}}  M_\alpha
\end{equation*}
is the Casimir element of $\mathrm{Hom}_{D_{t(\alpha)}}(M, D_{t(\alpha)}) \otimes_{D_{t(\alpha)}}  M_\alpha$. We also introduce the notation that $\overline{\alpha^*} = \{a^* \mid a\in \underline{\alpha} \}$ and $\underline{\alpha^*} = \{a^* \mid a\in \overline{\alpha} \}$.

\begin{mydef}\label{definition - S(S^1, S^2)}\cite[Definition 3.30]{soderberg2024mutation}
	Let $S^i$ be an acyclic species for each $i\in \{1, 2\}$. 
	\begin{enumerate}
		\item Define $S(S^1, S^2) = (D, M)$ as the species over $Q^1\tilde{\otimes}Q^2$ given by
		\begin{equation*}
			\begin{aligned}
				D &= D^1\otimes_\K D^2, \\
				M &= M^1\otimes_\K D^2\oplus D^1\otimes_\K M^2\oplus {M^1}^*\otimes_\K {M^2}^*.
			\end{aligned}
		\end{equation*}
		Using $\t_1$ and $\t_2$ we identify
		\begin{equation*}
			M^* = {M^1}^*\otimes_\K D^2 \oplus D^1\otimes_\K {M^2}^* \oplus M^1\otimes_\K M^2.
		\end{equation*}
		We define
		\begin{equation*}
			\b_{S(S^1, S^2)}: M\otimes_{D^1\otimes_\K D^2} M^* \oplus M^*\otimes_{D^1\otimes_\K D^2} M \to D^1\otimes_\K D^2,
		\end{equation*}
		via the morphism
		\begin{equation*}
			\t := \t_1\otimes \t_2: D = D^1\otimes_\K D^2 \to \K.
		\end{equation*}
		In other words, $\b_{S(S^1, S^2)}$ is defined such that
		\begin{equation*}
			\b_{S(S^1, S^2)}((a\otimes b)\otimes (c\otimes d)) = \begin{cases}
				\b_{S^1}(a\otimes c)\otimes bd, & \text{if } a\in M^1, b\in D^2, c\in {M^1}^*, d\in D^2 \\
				\b_{S^1}(a\otimes c)\otimes bd, & \text{if } a\in {M^1}^*, b\in  D^2, c\in M^1, d\in D^2 \\
				ac\otimes \b_{S^2}(b\otimes d), & \text{if } a\in D^1, b\in M^2, c\in D^1, d\in {M^2}^* \\
				ac\otimes \b_{S^2}(b\otimes d), & \text{if } a\in D^1, b\in {M^2}^*, c\in D^1, d\in M^2 \\
				\b_{S^1}(a\otimes c)\otimes \b_{S^2}(b\otimes d), & \text{if } a\in M^1, b\in M^2, c\in {M^1}^*, d\in {M^2}^* \\
				\b_{S^1}(a\otimes c)\otimes \b_{S^2}(b\otimes d), & \text{if } a\in {M^1}^*, b\in {M^2}^*, c\in M^1, d\in M^2, \\
				0, & \text{otherwise.}
			\end{cases}
		\end{equation*}
		\item The potential $W(S^1, S^2)\in T(S(S^1, S^2))$ is defined as
		\begin{equation}\label{eq - potential for tensor species}
			\begin{aligned}
				W(S^1, S^2) =& \sum_{\substack{(\alpha, \beta)\in Q_1^1\times Q_1^2 \\ (a, b')\in \underline{\alpha}\times \overline{\beta}}} (a\otimes e_{s(\beta)})\otimes(a^*\otimes {b'}^*)\otimes(e_{t(\alpha)}\otimes b') + \\
				&-\sum_{\substack{(\alpha, \beta)\in Q_1^1\times Q_1^2 \\ (a', b)\in \overline{\alpha}\times \underline{\beta}}} (e_{s(\alpha)}\otimes b)\otimes({a'}^*\otimes b^*)\otimes(a'\otimes e_{t(\beta)}) = \\
				=& W_1 - W_2.
			\end{aligned}
		\end{equation}
		It is indeed a potential, i.e. $W(S^1, S^2)\in \mathcal{Z}_{D^1\otimes_\K D^2}(T(S^1, S^2))$, since $c_\alpha\in \mathcal{Z}_{D^1}(T(\overline{S^1}))$ and $c_\beta\in \mathcal{Z}_{D^2}(T(\overline{S^2}))$ for any arrows $\alpha\in Q_1^1$ and $\beta\in Q_1^2$ which is a consequence of \eqref{eq - casimir elements equations ass. to b}.
	\end{enumerate}
\end{mydef}

\begin{myprop}\label{proposition - preprojective algebra of tensor species}\cite[Proposition 3.33]{soderberg2024mutation}
	Let $S^i$ be acyclic $\K$-species for $i\in \{1, 2\}$. Let $\Lambda = T(S^1)\otimes_\K T(S^2)$. Then $\gldim(\Lambda) = 2$ and $\widehat{\Pi}_3(\Lambda) \cong \mathcal{P}(S(S^1, S^2), W(S^1, S^2))$.
\end{myprop}

\begin{myrem}
	It is important to note that $S(S^1, S^2)$ is not a $\K$-species in general. This is due to the fact that if $D_1$ and $D_2$ are two division $\K$-algebras, we do not necessarily have that $D_1\otimes_\K D_2$ is a division $\K$-algebra (see \cite[Remark~3.31]{soderberg2024mutation} or \cite[Remark~11.6]{soderberg2022preprojective}). It may even happen that $\Lambda = T(S^1)\otimes_\K T(S^2)$ and $\widehat{\Pi}_3(\Lambda)$ are not basic. This will be problematic when we consider $2$-APR tilting in Section~\ref{Section - Preproj, 2apr and cutmutation} since $2$-APR tilting assumes that the algebra is basic.
\end{myrem}

Recall the following result from \cite{soderberg2024mutation}.

\begin{myprop}\cite[Proposition~3.32]{soderberg2024mutation}\label{proposition - product is division alg}
	Let $S^1$ and $S^2$ be two species. If $D^1_i\otimes_\K D^2_j$ is a division $\K$-algebra for all $i\in Q^1_0$ and $j\in Q^2_0$, then $S(S^1, S^2)$ is a $\K$-species. In particular, if $D^1 = \K^{|Q^1_0|}$ and $D^2 = \K^{|Q^2_0|}$, i.e. in the case when $T(S^1)$ and $T(S^2)$ are given as path algebras over $\K$, then $S(S^1, S^2)$ is a $\K$-species.
\end{myprop}

\begin{mycor}
	In addition to the assumptions in Proposition~\ref{proposition - product is division alg} assume that the quivers $Q^1$ and $Q^2$ of $S^1$ and $S^2$ respectively are trees. Then the set of all cuts of $(S(S^1, S^2), W(S^1, S^2))$ is transitive under successive cut-mutation.
\end{mycor}

In general $D^1_i\otimes_\K D^2_j$ is not necessarily a division $\K$-algebra. Thus we may have to alter quiver of the species $S(S^1, S^2)$ to obtain a $\K$-species. For this we need the following result.

\begin{myprop}\label{proposition - morita equivalent k-species}\cite[Proposition 3.24]{soderberg2024mutation}
	Let $(S, W)$ be a finite dimensional species with potential. There is a Morita equivalent $\K$-species with potential $(S', W')$. In other words, $T(S)\mathrm{-mod}\cong T(S')\mathrm{-mod}$, $\widehat{T}(S)\mathrm{-mod}\cong \widehat{T}(S')\mathrm{-mod}$ and $\mathcal{P}(S, W)\mathrm{-mod}\cong \mathcal{P}(S', W')\mathrm{-mod}$. 
\end{myprop}

\begin{myrem}\label{remark - morita equiv given by a idemp e}
	The Morita equivalence is achieved by choosing a certain idempotent $e\in D$ such that $\widehat{T}(S') = e\widehat{T}(S)e$ and $\mathcal{P}(S', W') = e\mathcal{P}(S, W)e$. In fact $(S', W')$ is a species with potential over a quiver with cycles $Q'$. Moreover, by \cite[Remark 3.6]{soderberg2024mutation} there is a quiver morphism $\pi: Q' \to Q$. Following the construction in \cite[Lemma 3.4]{soderberg2024mutation} one can check that if $c\in Q_2'$, then $\pi(c)\in Q_2$. Thus if $C\subseteq Q_1$ is a cut for $(S, W)$ then $\pi^{-1}(C)\subseteq Q_1'$ is a cut for $(S', W')$ which we call the corresponding cut. Note however that $(S', W')$ may admit cuts that are not of this form.
\end{myrem}

\begin{myrem}
	Proposition~\ref{proposition - morita equivalent k-species} says that $S(S^1, S^2)$ can be replaced with a $\K$-species $S'$ whenever $S^1$ and $S^2$ are finite dimensional species. Although we still end up with a $\K$-species, it is unclear whether the set of all cuts of the new quiver with cycles $Q'$ for $S'$ is transitive under successive cut-mutations. For example, if $Q^1\tilde{\otimes}Q^2$ is simply connected $Q'$ is not simply connected in general. To deal with this issue we next restrict our attention to a certain type of species.
\end{myrem}

Let $F = \K$ and $G$ be a finite dimensional division $\K$-algebra.

\begin{myrem}
	The assumptions on $F$ and $G$ are more general than those in \cite[Section 6]{soderberg2024mutation}, where $G$ is assumed to be a Galois field extension of $F$. That assumption was used in order to explicitly compute examples. Here, we only require that $G$ is a finite dimensional division $\K$-algebra since we prove more implicit statements.
\end{myrem}

\begin{myprop}\label{proposition - F-conn morita equiv}
	For $k\in \{1, 2\}$, let $S^k$ be a species over a tree such that $D_i^k, M_\alpha^k\in \{\K, G\}$ for all $i\in Q_0^k$ and $\alpha\in Q_1^k$. Then the species with potential $(S(S^1, S^2), W(S^1, S^2))$ is Morita equivalent to a $\K$-species with potential $(\tilde{S}, \tilde{W})$, defined below, over a quiver with cycles $\tilde{Q}$ which is covered and has enough cuts. Moreover,
	\begin{enumerate}
		\item if there are $i\in Q_0^1$ and $j\in Q_0^2$ such that $D_i^1 = D_j^2 = \K$, then $\tilde{Q}$ is simply connected,
		\item if $D_i^1 = D_j^2 = G$ for all $i\in Q_0^1$ and $j\in Q_0^2$, then $\tilde{Q}$ is a disjoint union of simply connected components.
	\end{enumerate}
	In both cases the set of all cuts of $(\tilde{S}, \tilde{W})$ is transitive under successive cut-mutation.
\end{myprop}

\begin{myrem}$\space$
	\begin{enumerate}
		\item Theorem~\ref{theorem - qwc fully, covered enough implies trans} applies to the quiver with cycles of $(\tilde{S}, \tilde{W})$ in Proposition~\ref{proposition - F-conn morita equiv}.
		\item Note that not all cases are covered in (1) and (2). Indeed, $(\tilde{S}, \tilde{W})$ may not be simply connected in the remaining cases. See Example~\ref{example - s tilde not simply connected}.
	\end{enumerate}
\end{myrem}

\begin{myrem}\label{remark - F = K}
	Note that Proposition~\ref{proposition - F-conn morita equiv} applies when $S^1$ and $S^2$ are representation finite species with $F = \K$.
\end{myrem}

Due to Remark~\ref{remark - F = K} we let $F = \K$ for the rest of the section. Before we prove Proposition~\ref{proposition - F-conn morita equiv}, let us first construct $\tilde{S}$ and $\tilde{W}$. We will do this in several steps. First we construct a species $S'$ and a potential $W'$ by following the proof of \cite[Lemma 3.4]{soderberg2024mutation} which will yield an isomorphic tensor algebra and Jacobian algebra. Then by following the construction in the proof of \cite[Lemma 3.24]{soderberg2024mutation} we get a Morita equivalent $\K$-species $\tilde{S}$ and a potential $\tilde{W}$ such that $T(S)\mathrm{-mod}\cong T(\tilde{S})\mathrm{-mod}$, $\widehat{T}(S)\mathrm{-mod}\cong \widehat{T}(\tilde{S})\mathrm{-mod}$ and $\mathcal{P}(S, W)\mathrm{-mod}\cong \mathcal{P}(\tilde{S}, \tilde{W})\mathrm{-mod}$. Before giving the general construction we give one example, which can be deduced from \cite{soderberg2024mutation}.

\begin{myex}\label{example - modify quiver}
	Let $S^1 = S^2$ be the $\R$-species $\C \xrightarrow{\C} \R$. Then $T(S^1)\otimes_\R T(S^2)$ is given by
	\begin{equation*}
		\begin{tikzcd}[column sep=1cm, row sep=1cm]
			\C\otimes_\R \C \ar[r, "\C\otimes_\R \C"] \ar[d, "\C\otimes_\R \C"] & \C\otimes_\R \R \ar[d, "\C\otimes_\R \R"] \\
			\R \otimes_\R \C \ar[r, "\R \otimes_\R \C"] & \R\otimes_\R \R
		\end{tikzcd}
	\end{equation*}
	or equivalently
	\begin{equation*}
		\begin{tikzcd}[column sep=1cm, row sep=1cm]
			\C\otimes_\R \C \ar[r, "\C\otimes_\R \C"] \ar[d, "\C\otimes_\R \C"] & \C \ar[d, "\C"] \\
			\C \ar[r, "\C"] & \R
		\end{tikzcd}
	\end{equation*}
	and certain relations. However $\C\otimes_\R \C$ is not a division $\R$-algebra, but $\C\otimes_\R \C \cong \C\oplus \C$, and thus we need to modify the quiver. According to \cite[Table 1]{soderberg2024mutation} $T(S^1)\otimes_\R T(S^2)$ is
	\begin{equation*}
		\begin{tikzcd}
			\C \arrow[rr, "\C"] \arrow[dd, "\C"] & & \C \arrow[dl, swap, "\C"] \\
			& \R \\
			\C \arrow[ur, swap, "\C"] & & \C \arrow[ll, "\C"] \arrow[uu, "\C"]
		\end{tikzcd}.
	\end{equation*}
\end{myex}

Let us now describe the general construction. Let $S^1 = (D^1, M^1)$ and $S^2 = (D^2, M^2)$ be species over $Q^1$ and $Q^2$ respectively. Set $S = S(S^1, S^2)$ and $W = W(S^1, S^2)$. Let us first describe $D = D^1\otimes_\K D^2$. Note that $F\otimes_\K F \cong F$, $F\otimes_\K G\cong G$ and $G\otimes_\K F\cong G$. For the last possibility $G\otimes_\K G$ we apply the Artin-Wedderburn theorem to get that
\begin{equation*}
	\phi: G\otimes_\K G \cong \bigoplus_{k=1}^nM_{N_k}(D_k'),
\end{equation*}
where $N_k$ and $n$ are positive integers and $D_k'$ are division $\K$-algebras. To simplify the notation we let $Q = Q^1\tilde{\otimes}Q^2$ and $V = \{(i, j)\in Q_0 \mid D_i^1 = D_j^2 = G\}$. We proceed to modify the quiver $Q$. Just as in Example~\ref{example - modify quiver}, the vertices in $V$ will be split up, whereas $Q_0\backslash V$ will remain. In the process arrows starting or ending in $V$ will be modified similarly. Let $x = (i, j)\in Q_0$. If $x\in V$ then we let
\begin{equation*}
	e_{xk} = \phi^{-1}(1_{M_{N_k}(D_k')}) \in D_i^1\otimes_\K D_j^2 \subseteq D^1\otimes_\K D^2.
\end{equation*}
We set
\begin{equation*}
	m_x = \begin{cases}
		n, & \text{if } x\in V \\
		1, & \text{if } x\in Q_0\backslash V
	\end{cases}
\end{equation*}
If $x \in Q_0 \backslash V$ we let $e_{xk} = 1_{D_i^1\otimes_\K D_j^2}$ so that
\begin{equation*}
	1_{D_i^1\otimes_\K D_j^2} = \sum_{k=1}^{m_x}e_{xk}
\end{equation*}
and
\begin{equation*}
	1_{T(S)} = 1_{D^1\otimes_\K D^2} = \sum_{x\in Q_0}\sum_{k=1}^{m_x}e_{xk}.
\end{equation*}
Let $Q'$ be the quiver given by
\begin{equation*}
	\begin{aligned}
		Q_0' &= \{(x, k) \mid x\in Q_0, k\in \{1, 2, \dots, m_x\} \}, \\
		Q_1' &= \{\alpha_{kk'}: (s(\alpha), k) \to (t(\alpha), k') \mid \alpha\in Q_1, k\in \{1, 2, \dots, m_{s(\alpha)}\}, k'\in \{1, 2, \dots, m_{t(\alpha)}\} \}.
	\end{aligned}
\end{equation*}
Define the species $S'$ over $Q'$ as
\begin{equation*}
	\begin{aligned}
		D_{(x, k)} = e_{xk}D^1\otimes_\K D^2 e_{xk}, \\
		M_{\alpha_{kk'}} = e_{t(\alpha)k'}M_\alpha e_{s(\alpha)k}.
	\end{aligned}
\end{equation*}
If $x\in V$ then we identify $D_{(x, k)}$ with $M_{N_k}(D_k')$ via $\phi$. If $x = (i, j)\in Q_0\backslash V$ then $D_{(x, k)} = D_i^1\otimes_\K D_j^2$. In this case we identify $D_{(x, k)}$ with $F$ whenever $D_i^1 = D_j^2 = F$ and $G$ otherwise.

Note that if $s(\alpha), t(\alpha) \in V$ then
\begin{equation*}
	M_{\alpha_{kk'}} = \begin{cases}
		\phi^{-1}(M_{N_k}(D_k')), & \text{if }k = k', \\
		0, & \text{if } k\neq k'.
	\end{cases}
\end{equation*}
Therefore $S'$ can be seen as a species over the quiver $\tilde{Q} = \tilde{Q}(S^1, S^2)$ defined by
\begin{equation*}
	\begin{aligned}
		\tilde{Q}_0 =& Q_0' \\
		\tilde{Q}_1 =& \{\alpha_{kk'}: (s(\alpha), k) \to (t(\alpha), k') \mid \alpha\in Q_1, s(\alpha)\not\in V \text{ or } t(\alpha)\not\in V, k\in \{1, 2, \dots, m_{s(\alpha)}\}, k'\in \{1, 2, \dots, m_{t(\alpha)}\} \} \cup \\
		&\cup \{\alpha_{kk}: (s(\alpha), k) \to (t(\alpha), k) \mid \alpha\in Q_1, s(\alpha), t(\alpha)\in V, k\in \{1, 2, \dots, n\} \}.
	\end{aligned}
\end{equation*}
In other words, $\tilde{Q}$ is obtained by removing all arrows $\alpha_{kk'}$ where $M_{\alpha_{kk'}} = 0$. We also let $\b_{S'}$ be induced by $\b_S$. By construction we immediately have that $T(S)\cong T(S')$ and $\widehat{T}(S)\cong \widehat{T}(S')$. Recall that $Q$ is a quiver with cycles where the cycles are induced by $W$. Let $W'$ to be the image of $W$ under the isomorphism $\widehat{T}(S)\cong \widehat{T}(S')$. Let us now compute $\tilde{Q}_2$. For this we need the following useful lemma. 

\begin{mylemma}\label{lemma - c_alpha x 1 casimir}
	Let $\alpha\in Q^1_1$ and $c_\alpha\in M_\alpha^1 \otimes_{D^1}{M_\alpha^1}^*$ be the Casimir element. Then 
	\begin{equation*}
		c_\alpha \otimes 1_{D^2_j} = \sum_{a\in \underline{\alpha}} (a\otimes 1_{D^2_j})\otimes (a^*\otimes 1_{D^2_j}) \in (M_\alpha^1\otimes_\K D^2_j)\otimes_D (M_\alpha^1\otimes_\K D^2_j)^*
	\end{equation*}
	is a Casimir element, for any $j\in Q^2_0$.
\end{mylemma}

\begin{proof}
	This directly follows from the definition of $\b$ for $S(S^1, S^2)$.
\end{proof}

Now using that
\begin{equation*}
	1_{D^1_{s(\alpha)}\otimes_\K D^2_j} = \sum_{k=1}^{m_{(s(\alpha), j)}} e_{(s(\alpha), j)k} 
\end{equation*}
we have that
\begin{equation*}
	c_\alpha \otimes 1_{D^2_j} = \sum_{\substack{a\in \underline{\alpha} \\ k\in \{1, 2, \dots, m_{(s(\alpha), j)}\}}} (a\otimes 1_{D^2_j})e_{(s(\alpha), j)k}\otimes e_{(s(\alpha), j)k} (a^*\otimes 1_{D^2_j}).
\end{equation*}

\begin{mylemma}\label{lemma - fixed k casimir}
	For a fixed $k$
	\begin{equation}\label{eq - idemp casimir element}
		\sum_{a\in \underline{\alpha}} (a\otimes 1_{D^2_j})e_{(s(\alpha), j)k}\otimes e_{(s(\alpha), j)k} (a^*\otimes 1_{D^2_j}) \in (M_\alpha^1\otimes_\K D^2_j)e_{(s(\alpha), j)k}\otimes_D ((M_\alpha^1\otimes_\K D^2_j)e_{(s(\alpha), j)k})^*
	\end{equation}
	is a Casimir element.
\end{mylemma}

\begin{proof}
	To prove this we need to verify the equations \eqref{eq - casimir elements equations ass. to b}. Let $x\in (M_\alpha^1\otimes_\K D^2_j)e_{(s(\alpha), j)k}$. Clearly, $xe_{(s(\alpha), j)k} = x$. Using that
	\begin{equation*}
		\begin{aligned}
			\b(e_{(s(\alpha), j)k} (a^*\otimes 1_{D^2_j}) \otimes x)\in D'_{((s(\alpha), j), k)} &= \t_r^{-1} ((a^*\otimes 1_{D^2_j}))(xe_{(s(\alpha), j)k}) \\
			&= \t_r^{-1} ((a^*\otimes 1_{D^2_j}))(x) = \\
			&= \b((a^*\otimes 1_{D^2_j}) \otimes x) \in D'_{((s(\alpha), j), k)}
		\end{aligned}
	\end{equation*}
	we get that
	\begin{equation*}
		\begin{aligned}
			\sum_{a\in \underline{\alpha}} (a\otimes 1_{D^2_j})e_{(s(\alpha), j)k}\otimes \b(e_{(s(\alpha), j)k} (a^*\otimes 1_{D^2_j}) \otimes x) &= \sum_{a\in \underline{\alpha}} (a\otimes 1_{D^2_j})e_{(s(\alpha), j)k}\otimes \b((a^*\otimes 1_{D^2_j}) \otimes x) = \\
			&= \sum_{a\in \underline{\alpha}} (a\otimes 1_{D^2_j})\otimes e_{(s(\alpha), j)k}\b((a^*\otimes 1_{D^2_j}) \otimes x) = \\
			&= \sum_{a\in \underline{\alpha}} (a\otimes 1_{D^2_j})\otimes \b((a^*\otimes 1_{D^2_j}) \otimes x) = x.
		\end{aligned}
	\end{equation*}
	The last equality follows due to Lemma~\ref{lemma - c_alpha x 1 casimir}. Similarly, the other equation in \eqref{eq - casimir elements equations ass. to b} is satisfied. Hence \eqref{eq - idemp casimir element} is a Casimir element.
\end{proof}

The new potential $W'$, i.e. the image of $W$, is explicitly described as
\begin{equation}\label{eq - potential W'}
	\begin{aligned}
		W' =& \sum_{\substack{(\alpha, \beta)\in Q_1^1\times Q_1^2 \\ (a, b')\in \underline{\alpha}\times \overline{\beta} \\ k\in \{1, 2, \dots, m_{(s(\alpha), s(\beta))}\} \\ k'\in \{1, 2, \dots, m_{(t(\alpha), t(\beta))}\} \\ k''\in \{1, 2, \dots, m_{(t(\alpha), s(\beta))}\} }} W_{a, b'}^+ - \sum_{\substack{(\alpha, \beta)\in Q_1^1\times Q_1^2 \\ (a', b)\in \overline{\alpha}\times \underline{\beta} \\ k\in \{1, 2, \dots, m_{(s(\alpha), s(\beta))}\} \\ k'\in \{1, 2, \dots, m_{(t(\alpha), t(\beta))}\} \\ k''\in \{1, 2, \dots, m_{(s(\alpha), t(\beta))}\} }} W_{a', b}^- \\
		W_{a, b'}^+ =& e_{(t(\alpha), s(\beta))k''}(a\otimes e_{s(\beta)})e_{(s(\alpha), s(\beta))k}\otimes e_{(s(\alpha), s(\beta))k}(a^*\otimes {b'}^*)e_{(t(\alpha), t(\beta))k'}\otimes e_{(t(\alpha), t(\beta))k'}(e_{t(\alpha)}\otimes b')e_{(t(\alpha), s(\beta))k''} \\
		W_{a', b}^- =& e_{(s(\alpha), t(\beta), k'')}(e_{s(\alpha)}\otimes b)e_{(s(\alpha), s(\beta))k}\otimes e_{(s(\alpha), s(\beta))k}({a'}^*\otimes b^*)e_{(t(\alpha), t(\beta))k'}\otimes e_{(t(\alpha), t(\beta))k'}(a'\otimes e_{t(\beta)})e_{(s(\alpha), t(\beta), k'')}
	\end{aligned}
\end{equation}
which is of the same form as \eqref{eq - potential for tensor species}. Thus
\begin{equation*}
	\tilde{Q}_2 = \{\alpha_{kk'}\beta_{k'k''}\gamma_{k''k} \mid \alpha_{kk'}, \beta_{k'k''}, \gamma_{k''k}\in \tilde{Q}_1, \alpha\beta\gamma\in Q_2, k\in \{1, \dots, m_{t(\alpha)}\}, k'\in \{1, \dots, m_{t(\beta)}\}, k''\in \{1, \dots, m_{t(\gamma)}\} \}.
\end{equation*}

Finally we construct the Morita equivalent species with potential $\tilde{S}$ and $\tilde{W}$ over the same quiver with cycles $\tilde{Q}$ in the following way. For every $(x, k)\in \tilde{Q}_0$ we let
\begin{equation*}
	\tilde{e}_{xk} = \phi^{-1}\left(\begin{bmatrix}
		1 & 0 & \cdots & 0 \\
		0 & 0 & & 0 \\
		\vdots & & \ddots & \vdots \\
		0 & 0 & \cdots & 0
	\end{bmatrix}\in M_{N_k}(D'_k)\right).
\end{equation*}
Now let $\tilde{S}(S^1, S^2)$ be the species over $\tilde{Q} = \tilde{Q}(S^1, S^2)$ defined as
\begin{equation*}
	\begin{aligned}
		\tilde{D}_{(x, k)} = \tilde{e}_{xk}D_{(x, k)}'\tilde{e}_{xk} \cong D_{k}' \\
		\tilde{M}_{\alpha_{kk'}} = \tilde{e}_{t(\alpha)k'}M_{\alpha_{kk'}}\tilde{e}_{s(\alpha)k}.
	\end{aligned}
\end{equation*}
and
\begin{equation*}
	\begin{aligned}
		\tilde{W} =& \sum_{\substack{(\alpha, \beta)\in Q_1^1\times Q_1^2 \\ (a, b')\in \underline{\alpha}\times \overline{\beta} \\ k\in \{1, 2, \dots, m_{(s(\alpha), s(\beta))}\} \\ k'\in \{1, 2, \dots, m_{(t(\alpha), t(\beta))}\} \\ k''\in \{1, 2, \dots, m_{(t(\alpha), s(\beta))}\} }} \tilde{W}_{a, b'}^+ - \sum_{\substack{(\alpha, \beta)\in Q_1^1\times Q_1^2 \\ (a', b)\in \overline{\alpha}\times \underline{\beta} \\ k\in \{1, 2, \dots, m_{(s(\alpha), s(\beta))}\} \\ k'\in \{1, 2, \dots, m_{(t(\alpha), t(\beta))}\} \\ k''\in \{1, 2, \dots, m_{(s(\alpha), t(\beta))}\} }} \tilde{W}_{a', b}^- \\
		\tilde{W}_{a, b'}^+ =& \tilde{e}_{(t(\alpha), s(\beta))k''}(a\otimes e_{s(\beta)})\tilde{e}_{(s(\alpha), s(\beta))k}\otimes \tilde{e}_{(s(\alpha), s(\beta))k}(a^*\otimes {b'}^*)\tilde{e}_{(t(\alpha), t(\beta))k'}\otimes \tilde{e}_{(t(\alpha), t(\beta))k'}(e_{t(\alpha)}\otimes b')\tilde{e}_{(t(\alpha), s(\beta))k''} \\
		\tilde{W}_{a', b}^- =& \tilde{e}_{(s(\alpha), t(\beta), k'')}(e_{s(\alpha)}\otimes b)\tilde{e}_{(s(\alpha), s(\beta))k}\otimes \tilde{e}_{(s(\alpha), s(\beta))k}({a'}^*\otimes b^*)\tilde{e}_{(t(\alpha), t(\beta))k'}\otimes \tilde{e}_{(t(\alpha), t(\beta))k'}(a'\otimes e_{t(\beta)})\tilde{e}_{(s(\alpha), t(\beta), k'')}
	\end{aligned}
\end{equation*}
Again, we let $\b_{\tilde{S}}$ to be induced by $\b_S$. To verify that the support of $\tilde{W}$ give rise to the cycles in $\tilde{Q}_2$, one needs to apply the same reasoning as in the proof of Lemma~\ref{lemma - fixed k casimir}.

\begin{proof}(Proof of Proposition~\ref{proposition - F-conn morita equiv})
	The cuts in $\tilde{Q}$ that correspond to the cuts $C_1$, $C_2$ and $C_3$ in Lemma~\ref{lemma - finite tress product trans} show that $\tilde{Q}$ is covered and has enough cuts.

	(1) It is left to show that $\tilde{Q}$ is simply connected. Let us first look at the quivers $\tilde{Q}(F, S^2)$ and $\tilde{Q}(G, S^2)$. Here $F$ and $G$ are viewed as species over a quiver with one vertex. Similarly we have the quivers $\tilde{Q}(S^1, F)$ and $\tilde{Q}(S^1, G)$. Now choose an inclusion $\iota: \tilde{Q}(F, S^2) \to \tilde{Q}(G, S^2)$ and a projection $\pi: \tilde{Q}(G, S^2) \to \tilde{Q}(F, S^2)$ such that $\pi\circ \iota = id_{\tilde{Q}(F, S^2)}$.

	It is enough to show that every cycle that lies entirely on the $1$-cells in $X_{\tilde{Q}}$ is homotopic to a point. First note that such a cycle is homotopic to a walk of the form
	\begin{equation*}
		\gamma_1\delta_1\gamma_2\delta_2\cdots \gamma_n\delta_n
	\end{equation*}
	where $\gamma_i$ are walks in $\bigcup_{j\in Q^2_0} \tilde{Q}(S^1, D^2_j)$ and $\delta_i$ are walks in $\bigcup_{i\in Q^1_0}\tilde{Q}(D^1_i, S^2)$. We will show that this walk is homotopic to a point by using induction on the number of vertices in $Q^1$. If $|Q^1_0| = 1$, then since $F$ appears in $S^1$ we have $S^1 = F$. Since $\tilde{Q}(F, S^2) \cong Q^2$ is a tree it is simply connected.

	For the induction step we choose a suitable orientation of $Q^1$. This leads to no loss of generality since the canvas does not depend on this orientation. More specifically, orient $Q^1$ so that there is an arrow $\alpha\in Q^1_1$ where $s(\alpha)$ is the only neighbour of $t(\alpha)$ and there is $k\in Q_0^1\backslash \{t(\alpha)\}$ where $D_k^1 = F$. This is possible since $Q^1$ is a tree and thus has two leaves. Note in particular, that if $Q^1$ has only one arrow the first condition is automatic. This ensures that $F$ still appears in the subspecies over the quiver where we have removed $t(\alpha)$ and $\alpha$ from $Q^1$. We have four cases for $D_{s(\alpha)}^1\xrightarrow{M_\alpha^1}D_{t(\alpha)}^1$.
	\begin{enumerate}[label=(\alph*)]
		\item $F\xrightarrow{M_\alpha^1}F$
		\item $F\xrightarrow{M_\alpha^1}G$
		\item $G\xrightarrow{M_\alpha^1}G$
		\item $G\xrightarrow{M_\alpha^1}F$
	\end{enumerate}
	For all these cases we will use the following argument. Pick a walk that goes through $X_{\tilde{Q}(D_{t(\alpha)}^1, S^2)}\subseteq X_{\tilde{Q}}$. Such a walk is of the form
	\begin{equation*}
		\beta_3(\alpha, v)_{dc}^{-1} \beta_2 (\alpha, u)_{ba}\beta_1,
	\end{equation*}
	for some walks $\beta_1$, $\beta_2$ and $\beta_3$, vertices $u,v\in Q_2$ and positive integers $a,b,c$ and $d$, where $\beta_2$ lies in $X_{\tilde{Q}(D^1_{t(\alpha)}, S^2)}$. If we can show that
	\begin{equation*}
		(\alpha, v)_{dc}^{-1} \beta_2 (\alpha, u)_{ba}
	\end{equation*}
	is homotopic to a walk that does not go through $X_{\tilde{Q}(D_{t(\alpha)}^1, S^2)}$ then we are done by induction. Via an inductive argument on the length of $\beta_2$, it is enough to show for the case when $\beta_2$ has length one. Again since the orientation of $Q^2$ does not matter, we can just pick an orientation such that $\beta_2 = (t(\alpha), \gamma)_{cb}$ where $\gamma\in Q^2_1$. Thus our walk is
	\begin{equation*}
		(\alpha, t(\gamma))^{-1}_{dc}(t(\alpha), \gamma)_{cb} (\alpha, s(\gamma))_{ba}.
	\end{equation*}

	The first three cases follow naturally since we have the following adjacent $2$-cells.
	\begin{equation*}
		\begin{tikzcd}[column sep=2cm, row sep=2cm]
			(s(\alpha), s(\gamma)) \arrow[r, "(\alpha{,}s(\gamma))_{ba}"] \arrow[d, "(s(\alpha){,}\gamma)_{da}"] & (t(\alpha), s(\gamma)) \arrow[d, "(t(\alpha){,} \gamma)_{cb}"] \\
			(s(\alpha), t(\gamma)) \arrow[r, "(\alpha{,} t(\gamma))_{cd}"] & (t(\alpha), t(\gamma)) \ar[ul, "(\alpha^*{,} \gamma^*)_{ad}", swap]
		\end{tikzcd}
	\end{equation*}
	In other words, we have the following homotopic maps
	\begin{equation*}
		(\alpha, t(\gamma))_{cd}^{-1} (t(\alpha), \gamma)_{cb} (\alpha, s(\gamma))_{ba} \sim (s(\alpha, \gamma))_{da}.
	\end{equation*}

	For the last case $G\xrightarrow{M_\alpha}F$ we need to work a little. We will prove this case in two steps. In the first step let $i, j\in Q_2$ such that $D_i^2 = F$. Let $\delta_j$ be a walk in $X_{\tilde{Q}(D_{t(\alpha)}^1, S^2)}$ from $(s(\alpha), i)$ to $(s(\alpha), j)$. We claim that
	\begin{equation*}
		(\alpha, j)_{1a} \sim \pi(\delta_j) (\alpha, i)_{11} \delta_j^{-1}.
	\end{equation*}
	Note that $\pi(\delta_j)\in X_{\tilde{Q}(D_{t(\alpha)}^1, S^2)}$ is defined as $D_{s(\alpha)}^1 = G$ and $D_{t(\alpha)}^1 = F$. As before we may assume that $\delta_j$ has length $1$ and the claim follows as in the previous cases.

	Now the result follows by
	\begin{equation*}
		\begin{aligned}
			(\alpha, t(\gamma))_{1d}^{-1} (t(\alpha), \gamma)_{11} (\alpha, s(\gamma))_{1a} &\sim \delta_{t(\gamma)} (\alpha, i)_{11}^{-1} \pi(\delta_{t(\gamma)})^{-1} (t(\alpha), \gamma)_{11} \pi(\delta_{s(\gamma)}) (\alpha, i)_{11} \delta_{s(\gamma)}^{-1} \sim \\
			&\sim \delta_{t(\gamma)} (\alpha, i)_{11}^{-1} (\alpha, i)_{11} \delta_{s(\gamma)}^{-1} \sim \\
			&\sim \delta_{t(\gamma)} \delta_{s(\gamma)}^{-1}
		\end{aligned}
	\end{equation*}
	where the second homotopy equivalence follows from $\tilde{Q}(D_{t(\alpha)}^1, S^2)$ being simply connected. This completes the proof of (1).

	(2) Now let us assume that $D_i^1 = D_j^2 = G$ for all $i\in Q_0^1$ and $j\in Q_0^2$. In other words, $S^1 = G \otimes_\K \K Q^1$ and $S^2= G \otimes_\K \K Q^2$. Note that
	\begin{equation*}
		T(S)\cong (G\otimes_\K G)\otimes_\K \K (Q^1\tilde{\otimes} Q^2) \cong \bigoplus_{k=1}^n M_{N_k}(D_k')\otimes_\K \K (Q^1\tilde{\otimes} Q^2).
	\end{equation*}
	Thus $T(\tilde{S}) = \bigoplus_{k=1}^n D_k'\otimes_\K \K (Q^1\tilde{\otimes} Q^2)$. The quiver with cycles $\tilde{Q}$ has $n$ connected components corresponding to $D_k'\otimes_\K \K (Q^1\tilde{\otimes} Q^2)$. All connected components are simply connected by the same argument as in the proof of Lemma~\ref{lemma - finite tress product trans}.
\end{proof}

\iffalse

\begin{myrem}\label{remark - rep finite F connected}
	In the more general case when we have field extensions $\K \subseteq F \subseteq G$ the quiver of $S'$ is not necessarily connected. Applying the argument in the proof of Proposition~\ref{proposition - F-conn morita equiv} implies that $Q'$ will be a quiver with cycles which is covered and has enough cuts such that each connected component is simply connected. In this more general setting, all representation finite species are species of type $(F \mid G)$ (see \cite{dlab1975algebras}). 
\end{myrem}

\fi

\begin{myex}\label{example - s tilde not simply connected}
	Let $\K = \R$, $F = \R$ and $G = \C$. Using \cite[Lemma 6.12]{soderberg2024mutation} we have that $\C \otimes \C \cong \C \times \C$. Let us first produce a non-example. Let $S^1 = G$ and $S^2 = F \xrightarrow{G}G\xrightarrow{G}F$. Then $T(S(S^1, S^2))\cong T(S')$ where $S'$ is the species
	\begin{equation*}
		\begin{tikzcd}[row sep=0cm]
			& G \arrow[dr, "G"] \\
			G \arrow[ru, "G"] \arrow[dr, "G"] & & G \\
			& G \arrow[ru, "G"]
		\end{tikzcd}.
	\end{equation*}
	Here $Q_2' = 0$ and therefore $X_{Q'}$ is homeomorphic to a circle. Hence $Q'$ is not simply connected.

	Let us modify the above example so it fits in Proposition~\ref{proposition - F-conn morita equiv}. If we let $S^1 = G\xrightarrow{G}F$. Then $T(S(S^1, S^2))\cong T(S')$ where $S'$ is the species
	\begin{equation*}
		\begin{tikzcd}[row sep=1cm, column sep=2cm]
			& G \arrow[rrd, "G"] \arrow[dddd, "G"] \arrow[from = ddddrr, "G"] \\
			G \arrow[ru, "G"] \arrow[ddd, "G"] \arrow[rr, "G\qquad", crossing over] & & G \arrow[r, "G"] \arrow[dddl, "G", crossing over] & G \arrow[ddd, "G"] \\ \\ \\
			G \arrow[r, "G"] & G \arrow[rr, "G"] \arrow[uuul, "G", swap] & & G \arrow[uuul, "G", swap]
		\end{tikzcd}.
	\end{equation*}
	Here $Q_2'$ consists of all cycles of order $3$. The space $X_{Q'}$ is homeomorphic to a $2$-dimensional disk and thus simply connected.
\end{myex}

\section{Truncated Jacobian Algebras}\label{section - truncated jacobian algebras}
In the spirit of Proposition~\ref{proposition - morita equivalent k-species} and Proposition~\ref{proposition - F-conn morita equiv} we assume from now on that all species are $\K$-species. In this section we define the truncated Jacobian algebras for species with potentials and cuts. We define the notion of algebraic cuts and prove that every cut of a self-injective species with potential is algebraic.

From the definition of the tensor algebra of a species $S$ we have a natural (completed) grading
\begin{equation*}
	\widehat{T}(S) = \prod_{i\ge 0}\widehat{T}(S)_i
\end{equation*}
where $\widehat{T}(S)_i = M^i$. For every subset $C\subseteq Q_1$, the grading $g_C$ induces a grading on $M$ by
\begin{equation*}
	M = M_0 \oplus M_1, \quad M_0=\bigoplus_{\alpha\not\in C}M_\alpha, \quad M_1 = \bigoplus_{\alpha\in C}M_\alpha.
\end{equation*}
We define $S_C = (D_i, M_\alpha)_{i\in Q_0, \alpha\in (Q_C)_1}$ over the quiver $Q_C = (Q_0, Q_1\backslash C)$. The grading induced by $g_C$ makes $\widehat{T}(S)$ into a bigraded $\K$-algebra, where
\begin{equation*}
	\begin{aligned}
		\widehat{T}(S)_{i, j} &= \bigoplus_{k_0 + k_1 + \dots + k_j = i - j} M_0^{k_0} \otimes_D M_1 \otimes_D M_0^{k_1} \otimes_D M_1 \otimes_D \cdots \otimes_D M_0^{k_{j-1}} \otimes_D M_1 \otimes_D M_0^{k_j}, \\
		\widehat{T}(S)_{i, \bullet} &= M^i, \\
		\widehat{T}(S)_{\bullet, j} &= \prod_{k_0, k_1, \dots, k_j\in \Z_{\ge 0}} M_0^{k_0} \otimes_D M_1 \otimes_D M_0^{k_1} \otimes_D M_1 \otimes_D \cdots \otimes_D M_0^{k_{j-1}} \otimes_D M_1 \otimes_D M_0^{k_j}.
	\end{aligned}
\end{equation*}
Note that $\widehat{T}(S_C) = \widehat{T}(S)_{\bullet, 0}$.

Now let $(S, W)$ be a species with potential and $C$ a cut. The grading induced by $C$ makes the Jacobian algebra $\mathcal{P}(S, W)$ a graded $\K$-algebra. More explicitly, let $R_\alpha = \{\partial_{\xi}(W) \mid \xi\in M_\alpha^*\}$. Then we can write
\begin{equation*}
	\begin{aligned}
		\mathcal{P}(S, W) &= \prod_{j\ge 0}\widehat{T}(S)_{\bullet, j} / \overline{I}_j \\
		I_j &= \sum_{\substack{\alpha\in C \\ x + y = j}}\widehat{T}(S)_{\bullet, x} R_\alpha \widehat{T}(S)_{\bullet, y} + \sum_{\substack{\alpha\not\in C \\ x + y = j-1}}\widehat{T}(S)_{\bullet, x} R_\alpha \widehat{T}(S)_{\bullet, y},
	\end{aligned}
\end{equation*}
where $\overline{I}_j$ is the closure of the $\widehat{T}(S_C)$-submodule $I_j\subseteq \widehat{T}(S)_{\bullet, j}$ with respect to the $\langle M_0 \rangle$-adic topology. We define the truncated Jacobian algebra $\mathcal{P}(S, W, C)$ as
\begin{equation*}
	\mathcal{P}(S, W, C) = \widehat{T}(S_C)/\mathcal{J}(S, W, C) = \widehat{T}(S_C)/\overline{\langle \partial_{a^*}(W) \mid a\in M_\alpha, \alpha\in C\rangle}.
\end{equation*}

\begin{myrem}\label{remark - Qc acyclic}
	In practice, $Q_C$ is an acyclic quiver since every known $2$-representation finite $\K$-algebra is acyclic. This is believed to be true in general, see \cite[Question 5.9]{HI14nrepinfinitealg}. In that case
	\begin{equation*}
		\widehat{T}(S)_{\bullet, j} = \bigoplus_{k_0, k_1, \dots, k_j =0}^N M_0^{k_0} \otimes_D M_1 \otimes_D M_0^{k_1} \otimes_D M_1 \otimes_D \cdots \otimes_D M_0^{k_{j-1}} \otimes_D M_1 \otimes_D M_0^{k_j},
	\end{equation*}
	where $N = \max\{k \mid M_0^k \not=0\}$. Thus $I_j = \overline{I}_j$ and we may write
	\begin{equation*}
		\mathcal{P}(S, W) = \prod_{j\ge 0}\widehat{T}(S)_{\bullet, j} / I_j.
	\end{equation*}
\end{myrem}

\begin{myprop}\label{proposition - degree zero of jacobian}
	For a species with potential $(S, W)$ and a cut $C$ we have that
	\begin{equation*}
		\mathcal{P}(S, W, C) = \mathcal{P}(S, W)_0
	\end{equation*}
	where the right hand side is the degree $0$ part with respect to the grading induced by $C$.
\end{myprop}

\begin{proof}
	From the description of $S_C$ we have that
	\begin{equation*}
		\widehat{T}(S)_{\bullet, 0} = \widehat{T}(S_C).
	\end{equation*}
	Also note that
	\begin{equation*}
		I_0 = \sum_{\substack{\alpha\in C}}\widehat{T}(S_C) R_\alpha \widehat{T}(S_C) = \langle \partial_{a^*}(W) \mid a^*\in M_\alpha^*, \alpha\in C\rangle
	\end{equation*}
	and therefore $\overline{I}_0 = \mathcal{J}(S, W, C)$.
\end{proof}

\begin{mydef}
	A cut $C\subseteq Q$ is called algebraic if $\dim_\K \mathcal{P}(S, W, C)<\infty$ and the sequence of left $\mathcal{P}(S, W, C)$-modules
	\begin{equation}\label{eq - alg cut}
		0\to \bigoplus_{\substack{\beta\in C \\ t(\beta) = i}}P_{s(\beta)}^{|\underline{\beta}|} \xrightarrow{(\partial_{\overline{\beta^*}, \underline{\alpha^*}}W)_{\alpha, \beta}} \bigoplus_{\substack{\alpha\not\in C \\ s(\alpha) = i}}P_{t(\alpha)}^{\overline{\alpha}}\xrightarrow{(\overline{\alpha})_\alpha} P_i\to S_i\to 0
	\end{equation}
	is exact, where $P_i = \mathcal{P}(S, W, C)e_i$.
\end{mydef}

\begin{myrem}$\space$
	\begin{enumerate}
		\item Note that if $C$ is an algebraic cut we immediately have that $\mathcal{P}(S, W, C)$ has global dimension at most $2$.
		\item The above definition of algebraic cuts differs from the definition given in \cite[Definition 3.2]{HerschendOsamu2011quiverwithpotential}. Our motivation is that we want \cite[Proposition 3.10]{HerschendOsamu2011quiverwithpotential} to hold in the species case, i.e. that every cut for a self-injective species with potential is algebraic. The minimality condition that they use in their definition is awkward to use in the species case since minimality depends on the $D$-$D$-bimodule structure of $M$.
	\end{enumerate}
\end{myrem}

\begin{myprop}\label{proposition - every cut is alg}
	If $(S, W)$ is a self-injective species with potential, then every cut $C\subseteq Q$ is algebraic. Moreover, every arrow in $Q$ appears at least once in $Q_2$.
\end{myprop}

\begin{proof}
	Let $\tilde{\Lambda} = \mathcal{P}(S, W)$. Since $(S, W)$ is self-injective, the sequence
	\begin{equation}\label{eq - alg cut proof}
		\tilde{P}_i\xrightarrow{(\underline{\beta})_\beta} \bigoplus_{\substack{\beta\in Q_1 \\ t(\beta) = i}}\tilde{P}_{s(\beta)}^{|\underline{\beta}|} \xrightarrow{(\partial_{\overline{\beta^*}, \underline{\alpha^*}}W)_{\alpha, \beta}} \bigoplus_{\substack{\alpha\in Q_1 \\ s(\alpha) = i}}\tilde{P}_{t(\alpha)}^{|\overline{\alpha}|}\xrightarrow{(\overline{\alpha})_\alpha} \tilde{P}_i\to S_i\to 0
	\end{equation}
	of $\tilde{\Lambda}$-modules is exact due to \cite[Proposition 5.2]{soderberg2024mutation}, where $\tilde{P}_i = \tilde{\Lambda}e_i$. By shifting degrees in \eqref{eq - alg cut proof} we get the exact sequence
	\begin{equation}\label{eq - alg cut proof degree}
		\begin{tikzcd}[column sep = 1.8cm]
			& \bigoplus_{\substack{\beta\in C \\ t(\beta) = i}}\tilde{P}_{s(\beta)}^{|\underline{\beta}|} \arrow[r, "(\partial_{\overline{\beta^*}, \underline{\alpha^*}}W)_{\alpha, \beta}"] & \bigoplus_{\substack{\alpha\not\in C \\ s(\alpha) = i}}\tilde{P}_{t(\alpha)}^{|\overline{\alpha}|} \arrow[r, "(\overline{\alpha})_\alpha"] & \tilde{P}_i \arrow[r] & 0 \\
			\tilde{P}_i(-1) \arrow[ru, "(\underline{\beta})_\beta"] \arrow[r, "(\underline{\delta})_\delta"] & \bigoplus_{\substack{\delta\not\in C \\ t(\delta) = i}}\tilde{P}_{s(\delta)}^{|\underline{\delta}|}(-1) \arrow[ru, "(\partial_{\overline{\delta^*}, \underline{\alpha^*}}W)_{\alpha, \delta}"] \arrow[r, "(\partial_{\overline{\delta^*}, \underline{\gamma^*}}W)_{\gamma, \delta}"] & \bigoplus_{\substack{\gamma\in C \\ s(\gamma) = i}}\tilde{P}_{t(\gamma)}^{|\overline{\gamma}|}(-1) \arrow[ru, "(\overline{\gamma})_\gamma"]
		\end{tikzcd}
	\end{equation}
	where the morphisms are of degree $0$. Now Proposition~\ref{proposition - degree zero of jacobian} implies that \eqref{eq - alg cut} is the degree $0$ part of \eqref{eq - alg cut proof degree}.

	For the second part note that if we assume that $\beta\in Q_1$ does not appear in $Q_2$, then it means that $\partial_{\overline{\beta^*}, \underline{\alpha^*} W} = 0$. This contradicts the exactness in \eqref{eq - alg cut proof} since the map $P_i \xrightarrow{\underline{\beta}}P_{s(\beta)}^{|\underline{\beta}|}$ is not surjective.
\end{proof}

\begin{myrem}\label{remark - jacobian as prerojecive}
	In the setting of Proposition~\ref{proposition - every cut is alg} we expect that $\mathcal{P}(S, W) = \Pi_3(\mathcal{P}(S, W, C))$. This is true for quivers with potentials, i.e. when $T(S) = \K Q$, by \cite{Keller2011CYcompletions}, see also \cite[Proposition 3.9]{HerschendOsamu2011quiverwithpotential}.
\end{myrem}

The following definition is motivated by Remark~\ref{remark - jacobian as prerojecive}.

\begin{mydef}
	Let $(S, W, C)$ be a species with potential and a cut $C$. We say that $C$ is a preprojective cut if $\gldim \mathcal{P}(S, W, C)\le 2$ and there is an isomorphism of complete graded algebras $\varphi: \widehat{\Pi}_3(\mathcal{P}(S, W, C)) \cong \mathcal{P}(S, W)$ with $\mathcal{P}(S, W)$ graded by $g_C$ such that $\varphi$ is the identity on the degree $0$ part.
\end{mydef}

\begin{myex}
	In Proposition~\ref{proposition - preprojective algebra of tensor species} we have that $\Lambda = T(S^1)\otimes_\K T(S^2)$ for acyclic species $S^i$. Let $S = S(S^1, S^2)$ then $\widehat{\Pi}_3(\Lambda) \cong \mathcal{P}(S, W)$. The cut $C_3$ in \eqref{eq - standard cuts} is a preprojective cut since $\mathcal{P}(S, W, C_3) \cong \Lambda$.

	By Proposition~\ref{proposition - morita equivalent k-species} there is a $\K$-species $\tilde{S}$ and a potential $\tilde{W}$ such that $T(S)\mathrm{-mod}\cong T(\tilde{S})\mathrm{-mod}$ and $\mathcal{P}(S, W)\mathrm{-mod}\cong \mathcal{P}(\tilde{S}, \tilde{W})\mathrm{-mod}$. By Remark~\ref{remark - morita equiv given by a idemp e} there is an idempotent $e\in \Lambda\subseteq \widehat{\Pi}_3(\Lambda)$ such that $e\widehat{\Pi}_3(\Lambda)e \cong \mathcal{P}(\tilde{S}, \tilde{W})$. Let $\tilde{C}_3$ be the corresponding cut induced from $C_3$ for $(\tilde{S}, \tilde{W})$. The cut $\tilde{C}_3$ induces a grading on $e\widehat{\Pi}_3(\Lambda)e$ such that 
	\begin{equation*}
		\mathcal{P}(\tilde{S}, \tilde{W}, \tilde{C}_3) = (e\widehat{\Pi}_3(\Lambda)e)_0 =: \Lambda'.
	\end{equation*}
	It follows that $\widehat{\Pi}_3(\Lambda') = \mathcal{P}(\tilde{S}, \tilde{W})$ is graded by $\tilde{C}_3$, so $\tilde{C}_3$ is also a preprojective cut.
\end{myex}

\begin{myquestion}
	Is every algebraic cut a preprojective cut? This is true in the setting of \cite{HerschendOsamu2011quiverwithpotential} due to \cite{Keller2011CYcompletions}, see also \cite{HerschendOsamu2011quiverwithpotential}. We expect that it is also true in our setting.
\end{myquestion}

\section{Preprojective Algebra, 2-APR tilts and Cut-Mutation}\label{Section - Preproj, 2apr and cutmutation}
In this section we first recall the definition and properties of $2$-APR tilts. Then we show that under certain conditions $2$-APR tilting of $\mathcal{P}(S, W, C)$ corresponds to cut-mutation on $C$.

\subsection{2-APR Tilting}
\begin{mydef}\cite[Definition 3.1]{IO11nRFalgandnAPRtilt}\label{definition - d-apr tilting module}
	Let $\Lambda$ be a basic $\K$-algebra and $P$ be a simple projective module such that $\Hom_\Lambda(\DD\Lambda, P) = \mathrm{Ext}_\Lambda^1(\DD\Lambda, P) = 0$. We decompose $\Lambda = P\oplus P'$ and call
	\begin{equation*}
		T = \tau_2^-(P)\oplus P'
	\end{equation*}
	the weak $2$-APR tilting $\Lambda$-module associated with $P$. Moreover, if $\mathrm{idim}(P) = 2$, then $T$ is a $2$-APR tilting $\Lambda$-module and $\mathrm{End}_\Lambda(T)^{op}$ is called a $2$-APR tilt of $\Lambda$.
\end{mydef}

\begin{myprop}\cite[Proposition 3.6]{IO11nRFalgandnAPRtilt}
	If $\mathrm{gl.dim}(\Lambda) = 2$ and $T$ is a $2$-APR tilting $\Lambda$-module, then $\mathrm{gl.dim}(\mathrm{End}_\Lambda(T))^{op} = 2$.
\end{myprop}

Note that if $\Lambda$ is $2$-representation finite then any simple projective non-injective $\Lambda$-module $P$ admits a $2$-APR tilt of $\Lambda$.

\begin{mycor}\cite[Corollary 4.3]{IO11nRFalgandnAPRtilt}
	Iterated $2$-APR tilts of a $2$-representation finite $\K$-algebra are $2$-representation finite.
\end{mycor}

\subsection{Preprojective Algebra and 2-APR tilts}
Let $\Gamma = \Gamma_\bullet = \prod_{i\ge 0} \Gamma_i$ be a complete graded $\K$-algebra such that $\Gamma_0 = \Lambda$ is a basic finite dimensional algebra with $\gldim(\Lambda) = 2$ and $\widehat{\Pi}_3(\Lambda) \cong \Gamma_\bullet$, i.e. $\Gamma_i \cong \Hom_{\mathcal{D}^b(\Lambda)}(\Lambda, \nu_2^{-i}(\Lambda))$. Let $e\in \Lambda$ be an idempotent such that $\Lambda e$ is simple and projective. Assume that $T = \tau_2^-(\Lambda e)\oplus \Lambda(1- e)$ is a $2$-APR tilting module. Let us first show that $\tau_2^-(\Lambda e)\cong \nu_2^{-1}(\Lambda e)$ in $\mathcal{D}^b(\Lambda)$. Let
\begin{equation*}
	0\to \Lambda e\to I_0\to I_1\to I_2 \to 0
\end{equation*}
be an injective resolution of $\Lambda e$. Then
\begin{equation*}
	\DD(I_2)\to \DD(I_{1}) \to \DD(I_0)\to \DD(\Lambda e)\to 0
\end{equation*}
is a projective resolution of $\DD(\Lambda e)$, where $\DD = \Hom_\K(-, \K)$. Now applying $\Hom_{\Lambda^{op}}(-, \Lambda)$ yields the complex
\begin{equation*}
	0\to \nu^{-1}(I_0)\to \nu^{-1}(I_1)\to \nu^{-1}(I_2)\to 0
\end{equation*}
where $\tau_2^-(\Lambda e)$ is given as the homology at $\nu^{-1}(I_2)$ in the above complex. By definition of $2$-APR tilting we have that $\mathrm{Ext}_\Lambda^i(\DD\Lambda, \Lambda e) = 0$ for $0\le i\le 1$. Hence we only have non-zero homology at $\nu^{-1}(I_2)$ and therefore $\nu_2^{-1}(\Lambda e)\cong \tau_2^-(\Lambda e)$.

Now let $\Lambda' = \mathrm{End}_\Lambda(T)^{op} = \mathrm{End}_{\mathcal{D}^b(\Lambda)}(\nu_2^{-1}(\Lambda e)\oplus \Lambda(1-e))^{op}$. Then by \cite[Proposition 3.6]{IO11nRFalgandnAPRtilt} $\gldim(\Lambda') = 2$. Let us now study $\widehat{\Pi}_3(\Lambda')$. First we write
\begin{equation*}
	\Gamma = \begin{bmatrix}
		e\Gamma e & e\Gamma (1-e) \\
		(1- e)\Gamma e & (1- e)\Gamma (1-e)
	\end{bmatrix}, \qquad \Gamma_i = \begin{bmatrix}
		e\Gamma_i e & e\Gamma_i (1-e) \\
		(1- e)\Gamma_i e & (1- e)\Gamma_i (1-e)
	\end{bmatrix}.
\end{equation*}
We define a new grading
\begin{equation*}
	\Gamma^i = \begin{bmatrix}
		e\Gamma_i e & e\Gamma_{i-1} (1-e) \\
		(1- e)\Gamma_{i+1} e & (1- e)\Gamma_i (1-e)
	\end{bmatrix}.
\end{equation*}
Note that since $\Lambda e$ is a simple projective we have that $(1- e)\Gamma_0 e = 0$, which implies that $\widehat{\Pi}_3(\Lambda)\cong \Gamma^\bullet = \prod_{i\ge 0}\Gamma^i$.

\begin{myprop}\label{proposition - degree shift for preprojective}
	There is an isomorphism $\widehat{\Pi}_3(\Lambda')\cong \Gamma^\bullet$ as complete graded algebras. In particular,
	\begin{equation*}
		\Lambda' \cong \Gamma^0 = \begin{bmatrix}
			e\Gamma_0 e & 0 \\
			(1- e)\Gamma_1 e & (1-e)\Gamma_0 (1-e)
		\end{bmatrix}.
	\end{equation*}
\end{myprop}

\begin{proof}
	Since $T$ is a tilting module we have an equivalence $\mathcal{D}^b(\Lambda')\xrightarrow{\sim}\mathcal{D}^b(\Lambda)$ by \cite[Theorem 1.6]{Happel1987}, that sends $\Lambda'$ to $T$. Thus
	\begin{equation*}
		\begin{aligned}
			\widehat{\Pi}_3(\Lambda') &= \prod_{i\ge 0} \Hom_{\mathcal{D}^b(\Lambda')}(\Lambda', \nu_2^{-i}(\Lambda')) = \\
			&= \prod_{i\ge 0} \Hom_{\mathcal{D}^b(\Lambda)}(T, \nu_2^{-i}(T)) = \\
			&= \prod_{i\ge 0} \Hom_{\mathcal{D}^b(\Lambda)}(\nu_2^{-1}(\Lambda e)\oplus \Lambda(1- e), \nu_2^{-i}(\nu_2^{-1}(\Lambda e)\oplus \Lambda(1- e))) = \\
			&= \prod_{i\ge 0} \Hom_{\mathcal{D}^b(\Lambda)}(\nu_2^{-1}(\Lambda e), \nu_2^{-i}(\nu_2^{-1}(\Lambda e))) \oplus \Hom_{\mathcal{D}^b(\Lambda)}(\nu_2^{-1}(\Lambda e), \nu_2^{-i}(\Lambda(1- e))) \\
			&\qquad \oplus \Hom_{\mathcal{D}^b(\Lambda)}(\Lambda(1- e), \nu_2^{-i}(\nu_2^{-1}(\Lambda e))) \oplus \Hom_{\mathcal{D}^b(\Lambda)}(\Lambda(1- e), \nu_2^{-i}(\Lambda(1- e))) \\
			&= \prod_{i\ge 0} \Hom_{\mathcal{D}^b(\Lambda)}(\Lambda e, \nu_2^{-i}(\Lambda e)) \oplus \Hom_{\mathcal{D}^b(\Lambda)}(\Lambda e, \nu_2^{-(i-1)}(\Lambda(1- e))) \\
			&\qquad \oplus \Hom_{\mathcal{D}^b(\Lambda)}(\Lambda(1- e), \nu_2^{-(i+1)}(\Lambda e)) \oplus \Hom_{\mathcal{D}^b(\Lambda)}(\Lambda(1- e), \nu_2^{-i}(\Lambda(1- e))) = \\
			&= \prod_{i\ge 0}\Gamma^i.
		\end{aligned}
	\end{equation*}
\end{proof}

\iffalse
\begin{mycor}
	Let $(S, W, C)$ be a species with potential with a cut $C$ such that $Q_C$ is acyclic, $\gldim(\mathcal{P}(S, W, C)) = 2$ and $\mathcal{P}(S, W)\cong \widehat{\Pi}_3(\mathcal{P}(S, W, C))$ as complete graded algebras, where the grading on $\mathcal{P}(S, W)$ is induced by $g_C$. If $k\in Q_0$ is a strict sink and gives a $2$-APR tilting module $T$, then $\Lambda' = \mathrm{End}_{\mathcal{P}(S, W, C)}(T)^{op}$ satisfies $\Lambda' \cong \mathcal{P}(S, W, \mu_k^-(C))$ and $\widehat{\Pi}_3(\Lambda')\cong \mathcal{P}(S, W)$ as complete graded algebras, where $\mathcal{P}(S, W)$ is graded via $g_{\mu_k^-(C)}$.
\end{mycor}

\begin{myrem}
	In \cite{HI14nrepinfinitealg} they conjectured that any $d$-representation finite $\K$-algebra enjoys some acyclic property. This means that, under the conjecture, that if $\mathcal{P}(S, W, C)$ is a $2$-representation finite $\K$-algebra then $Q_C$ is an acyclic quiver. Thus we have an easier description of $\mathcal{P}(S, W, C)$ due to Remark~\ref{remark - Qc acyclic}.
\end{myrem}

\begin{myrem}
	If $\mathcal{P}(S, W, C)$ is a $2$-representation finite $\K$-algebra, then every strict sink gives a $2$-APR tilting module $T$.
\end{myrem}

\begin{mycor}\label{corollary - cut mutation and 2-apr tilts}
	Let $(S, W)$ be a species with potential and $C$ be a cut, such that $\gldim(\mathcal{P}(S, W, C)) \le 2$ and is its $\mathcal{P}(S, W)$ is the $3$-preprojective algebra. Let $T$ be a $2$-APR tilting $\mathcal{P}(S, W, C)$-module associated to a sink $k$ in $Q_C$. Then
	\begin{equation}\label{eq - cut and apr equation}
		\mathrm{End}_{\mathcal{P}(S, W, C)}(T) \cong \mathcal{P}(S, W, \mu_k^-(C)).
	\end{equation}
\end{mycor}

\begin{proof}
	First note that if $\Lambda = \mathcal{P}(S, W, C)$, then $\mathrm{End}_{\mathcal{P}(S, W, C)}(T) = \Lambda'$ as in Theorem~\ref{proposition - degree shift for preprojective}. The isomorphism in \eqref{eq - cut and apr equation} is constructed by comparing $\mathcal{P}(S, W, \mu_k^-(C))$ and $\Lambda'$ and applying Theorem~\ref{proposition - degree shift for preprojective}.
\end{proof}

\begin{myrem}
	By Corollary~\ref{corollary - cut mutation and 2-apr tilts} we have a connection between cut-mutation and $2$-APR tilts. In this remark we will demonstrate why we assume that our quivers does not have multiple arrows. Let $S$ be the species such that $T(S) = \R Q$, where $Q$ is given as
	\begin{equation*}
		\begin{tikzcd}
			1 \arrow[bend left, r, "\alpha_1"] \arrow[bend right, r, "\alpha_2"] & 2 \arrow[d, "\beta"]\\
			3 \arrow[u, "\delta"] & 4 \arrow[l, "\gamma"]
		\end{tikzcd}
	\end{equation*}
	and a potential
	\begin{equation*}
		W = \delta \gamma \beta \alpha_1.
	\end{equation*}
	We have the following cuts of $Q$ which we denote by $C$ and $C'$ respectively.
	\begin{equation*}
		\begin{tikzcd}
			1 \arrow[bend left, r] \arrow[bend right, r] & 2 \arrow[d, dashed] & 1 \arrow[bend left, r, dashed] \arrow[bend right, r, dashed] & 2 \arrow[d] \\
			3 \arrow[u] & 4 \arrow[l] & 3 \arrow[u] & 4 \arrow[l]
		\end{tikzcd}
	\end{equation*}
	Clearly, we get the other cut by mutating at $2$. Also note that $C'$ is not an algebraic cut. One quickly checks that $\mathcal{P}(S, W, C)$ and $\mathcal{P}(S, W, C')$ are of global dimension $2$. Since $\partial_{\alpha_2^*}(W) = 0$ the $2$-APR tilt associated to the vertex $2$ of $\mathcal{P}(S, W, C')$ is isomorphic to the path algebra
	\begin{equation*}
		\begin{tikzcd}
			1 \arrow[r, dashed, -] & 2 \arrow[d, "\beta"] \\
			3 \arrow[u, "\delta"] & 4 \arrow[l, "\gamma"]
		\end{tikzcd}
	\end{equation*}
	with the relation $\alpha\delta\gamma = 0$.
\end{myrem}

\fi

\subsection{Cut-Mutation and 2-APR Tilts}
Recall the cut-mutation defined in Definition~\ref{definition - cut-mutation}. The following theorem relates cut-mutation and $2$-APR tilting for species with potentials.

\begin{mythm}\label{theorem - cut and apr correspondence}
	Let $(S, W)$ be a species with potential and let $C$ be a preprojective cut. If $k\in Q_0$ is a strict sink and gives a $2$-APR tilting module $T$, then $\mu_k^-(C)$ is a preprojective cut and
	\begin{equation}\label{eq - cut and apr equation}
		\mathrm{End}_{\mathcal{P}(S, W, C)}(T)^{op} \cong \mathcal{P}(S, W, \mu_k^-(C)).
	\end{equation}
\end{mythm}

\begin{proof}
	Let $\Lambda = \mathcal{P}(S, W, C)$ and $\Gamma_\bullet = \mathcal{P}(S, W)_\bullet$ be graded by $C$. Since $C$ is preprojective we get $\Gamma_\bullet \cong \widehat{\Pi}_3(\Lambda)$ which restricts to the identity on $\Gamma_0 = \Lambda$. The $2$-APR tilting module is given as $T = \tau_2^-(\Lambda e_k)\oplus \Lambda(1- e_k)$. By Proposition~\ref{proposition - degree shift for preprojective} $\Gamma^\bullet \cong \widehat{\Pi}_3(\Lambda')$ where $\Lambda' = \mathrm{End}_{\mathcal{P}(S, W, C)}(T)^{op}$. We are done if we can prove that $\Gamma^\bullet = \mathcal{P}(S, W)^\bullet$ graded by $\mu_k^-(C)$. Since $k$ is a strict sink in $Q_C$, $k$ is a strict source in $Q_{\mu_k^-(C)}$. Hence
	\begin{equation}\label{eq - cut and apr eq 1}
		e_k\mathcal{P}(S, W)_ie_k = e_k\mathcal{P}(S, W)^ie_k
	\end{equation}
	and
	\begin{equation}\label{eq - cut and apr eq 2}
		(1-e_k)\mathcal{P}(S, W)_i(1- e_k) = (1-e_k)\mathcal{P}(S, W)^i(1- e_k).
	\end{equation}
	Let $M_C = \bigoplus_{\alpha\in C}M_\alpha$ and $M_{\mu_k^-(C)} = \bigoplus_{\alpha\in \mu_k^-(C)}M_\alpha$. Then $(1- e_k)\Gamma e_k = (1-e_k)\Gamma M_C e_k$ and $e_k \Gamma(1- e_k) = e_k M_{\mu_k^-(C)}\Gamma(1-e_k)$. Hence
	\begin{equation*}
		(1-e_k)\Gamma_{i+1}e_k = (1-e_k)\mathcal{P}(S, W)_iM_C e_k.
	\end{equation*}
	Note that $M_C e_k = (1-e_k)M_C e_k$ and thus using \eqref{eq - cut and apr eq 2} we have
	\begin{equation}\label{eq - cut and apr eq 3}
		\begin{aligned}
			(1-e_k)\Gamma_{i+1}e_k &= (1-e_k)\mathcal{P}(S, W)_i(1-e_k)M_C e_k = (1-e_k)\mathcal{P}(S, W)^i(1- e_k)M_C e_k = \\
			&= (1-e_k)\mathcal{P}(S, W)^i e_k.
		\end{aligned}
	\end{equation}
	Similarly
	\begin{equation}\label{eq - cut and apr eq 4}
		\begin{aligned}
			e_k \Gamma_{i-1}(1-e_k) &= e_k M_{\mu_k^-(C)} \mathcal{P}(S, W)_{i-1} (1-e_k) = e_k M_{\mu_k^-(C)}(1-e_k) \mathcal{P}(S, W)^{i-1}(1-e_k) = \\
			&= e_k \mathcal{P}(S, W)^{i}(1-e_k).
		\end{aligned}
	\end{equation}
	Combining \eqref{eq - cut and apr eq 1}, \eqref{eq - cut and apr eq 2}, \eqref{eq - cut and apr eq 3} and \eqref{eq - cut and apr eq 4} gives $\Gamma^\bullet = \mathcal{P}(S, W)^\bullet$.
\end{proof}

Now we have enough to state the main theorem of the article.

\begin{mythm}\label{theorem - jacobian 2apr tils of each other}
	Let $(S, W)$ be a self-injective simply connected species with potential with enough cuts. Assume that $(S, W)$ has a preprojective cut. Then all cuts $C$ are preprojective and the corresponding truncated Jacobian algebras $\mathcal{P}(S, W, C)$ are iterated $2$-APR tilts of each other. In particular, each $\mathcal{P}(S, W, C)$ is $2$-representation finite.
\end{mythm}

\begin{proof}
	Let $C$ be a preprojective cut. Then $\mathcal{P}(S, W) \cong \widehat{\Pi}_3(\mathcal{P}(S, W, C))$ which is self-injective by assumption. Applying \cite[Corollary 3.7]{IO13Stablecateofhigherpreproj} gives that $\mathcal{P}(S, W, C)$ is a $2$-representation finite algebra. By Proposition~\ref{proposition - every cut is alg} every arrow in $Q$ appears at least once in $Q_2$ and thus $Q$ is covered. By applying Theorem~\ref{theorem - qwc fully, covered enough implies trans} we have that all cuts are transitive under successive cut-mutation. Thus using Theorem~\ref{theorem - cut and apr correspondence}, all truncated Jacobian algebras are iterated $2$-APR tilts of each other. By \cite[Corollary 4.3]{IO11nRFalgandnAPRtilt} $2$-APR tilting preserves $2$-representation finiteness, hence all truncated Jacobian algebras are $2$-representation finite.
\end{proof}

\begin{mycor}
	Let $S^1$ and $S^2$ be two $l$-homogeneous representation finite species such that $D^1_i\otimes_\K D^2_j$ is a division $\K$-algebra for all $i\in Q^1_0$ and $j\in Q^2_0$, then all truncated Jacobian algebras of $(S(S^1, S^2), W(S^1, S^2))$ are iterated $2$-APR tilts of each other.
\end{mycor}

\begin{proof}
	This follows by combining Lemma~\ref{lemma - finite tress product trans}, Proposition~\ref{proposition - product is division alg} and Theorem~\ref{theorem - jacobian 2apr tils of each other}.
\end{proof}

Similarly we have the following corollary by applying Proposition~\ref{proposition - F-conn morita equiv}.

\begin{mycor}\label{corollary - tensor product species morita transitive}
	Let $S^1$ and $S^2$ be two $l$-homogeneous representation finite species satisfying the conditions in Proposition~\ref{proposition - F-conn morita equiv} (1) or (2), then all truncated Jacobian algebras of $(\tilde{S}, \tilde{W})$ are iterated $2$-APR tilts of each other, where $(\tilde{S}, \tilde{W})$ is the Morita equivalent $\K$-species with potential given in Proposition~\ref{proposition - F-conn morita equiv}.
\end{mycor}

Consider the species in Figure~\ref{Figure - Dynkin Diagrams} with $F = \K$. They are all representation finite and Corollary~\ref{Corollary - l homogeneous species} determines which are $l$-homogeneous. Note that Corollary~\ref{corollary - tensor product species morita transitive} applies to any pair of them.

\begin{myex}
	Let us now compute an example where Theorem~\ref{theorem - jacobian 2apr tils of each other} can be applied. For this we use the same example constructed in Example~\ref{example - modify quiver}. Let $S$ be $\R$-species
	\begin{equation*}
		\C \xrightarrow{\C} \R.
	\end{equation*}
	Now $S$ is a $2$-homogeneous representation finite species, since $S$ is of type $B_2$, the algebra $T(S)\otimes_\R T(S)$ is a $2$-representation finite algebra (see \cite[Corollary 3.7]{soderberg2022preprojective} and \cite[Corollary 1.5]{herschend2011n}). According to the table \cite[Table 1]{soderberg2024mutation} the algebra $T(S)\otimes_\R T(S)$ is isomorphic to $\mathcal{P}(S', W, C)$, where $S'$ is the species
	\begin{equation*}
		\begin{tikzcd}
			\C_1 \arrow[rr, "\lsub{2}\C_1"] \arrow[dd, "\lsub{4}\C_1"] & & \C_2 \arrow[dl, swap, "\lsub{3}\C_2"] \\
			& \R_3 \arrow[ul, swap, "\lsub{1}\C_3"] \arrow[dr, swap, "\lsub{5}\C_3"] \\
			\C_4 \arrow[ur, swap, "\lsub{3}\C_4"] & & \C_5 \arrow[ll, "\lsub{4}\C_5"] \arrow[uu, "\lsub{2}\C_5"]
		\end{tikzcd}
	\end{equation*}
	over the quiver
	\begin{equation*}
		\begin{tikzcd}
			1 \arrow[rr, "\lsub{2}\alpha_1"] \arrow[dd, "\lsub{4}\alpha_1"] & & 2 \arrow[dl, swap, "\lsub{3}\alpha_2"] \\
			& 3 \arrow[ul, swap, "\lsub{1}\alpha_3"] \arrow[dr, swap, "\lsub{5}\alpha_3"] \\
			4 \arrow[ur, swap, "\lsub{3}\alpha_4"] & & 5 \arrow[ll, "\lsub{4}\alpha_5"] \arrow[uu, "\lsub{2}\alpha_5"]
		\end{tikzcd}
	\end{equation*}
	with the potential
	\begin{equation*}
		\begin{aligned}
			W =& \lsub{2}1_1\otimes \lsub{1}1^*_3\otimes \lsub{3}1_2 + \lsub{2}1_1\otimes \lsub{1}i^*_3\otimes \lsub{3}1_2 - \lsub{4}1_1\otimes \lsub{1}1^*_3\otimes \lsub{3}1_4 - \lsub{4}1_1\otimes \lsub{1}i^*_3\otimes \lsub{3}1_4 + \\
			&+ \lsub{2}1_5\otimes \lsub{5}1^*_3\otimes \lsub{3}1_2 + \lsub{2}1_5\otimes \lsub{5}i^*_3\otimes \lsub{3}1_2 - \lsub{4}1_5\otimes \lsub{5}1^*_3\otimes \lsub{3}1_4 + \lsub{4}1_5\otimes \lsub{5}i^*_3\otimes \lsub{3}1_4
		\end{aligned}
	\end{equation*}
	and the cut $C = \{\lsub{1}\alpha_3, \lsub{5}\alpha_3 \}$. By Proposition~\ref{proposition - F-conn morita equiv} $(S', W')$ is a simply connected species with potential, which is covered and has enough cuts. We display all cuts in Figure~\ref{figure - theorem 7.6 applied} and connected cuts via cut-mutation are noted with an edge.
	\begin{figure}[!ht]
	\begin{adjustbox}{width=\columnwidth,center}
		\begin{tikzcd}[bezier bounding box]
			& \Image{exthm76/3/3.pdf} \ar[thick, -, dr] & & \Image{exthm76/6/6.pdf} \ar[thick, -, dr] \\
			\Image{exthm76/1/1.pdf} \ar[thick, -, dr] \ar[thick, -, ur] & & \Image{exthm76/4/4.pdf} \ar[thick, -, dr] \ar[thick, -, ur] & & \Image{exthm76/7/7.pdf} \ar[llll, thick, -, in=90, out=90, looseness=1.2] \\
			& \Image{exthm76/2/2.pdf} \ar[thick, -, ur] & & \Image{exthm76/4/4.pdf} \ar[thick, -, ur]
		\end{tikzcd}
	\end{adjustbox}
		\caption{Mutation lattice for $B_2\times B_2$}\label{figure - theorem 7.6 applied}
	\end{figure}
\end{myex}

\subsection{Tensor Products of Representation Finite Species}
In this section we compute more examples where Corollary~\ref{corollary - tensor product species morita transitive} is applied.

\begin{myex}\label{ex - example for intro}
	Let $S^1$ be the species $\R \xleftarrow{\R}\R \xrightarrow{\R}\R$ and $S^2$ be the species $\C\xrightarrow{\C}\R$. Following \cite[Example 6.11]{soderberg2024mutation} the species $S(S^1, S^2)$ and $W(S^1, S^2)$ are given as the following.
	\begin{equation*}
		\begin{tikzcd}[row sep = 0.3cm, column sep = 1cm]
			& \C \arrow[r, "\lsub{2}\C_1"] \arrow[dd, "\lsub{4}\C_1"] & \C \arrow[dd, "\lsub{5}\C_2"] & \C \arrow[l, "\lsub{2}\C_3"'] \arrow[dd, "\lsub{6}\C_3"] \\
			S(S^1, S^2): \\
			& \R \arrow[r, "\lsub{5}\R_4"] & \R \arrow[luu, "\lsub{1}\C_5"'] \arrow[ruu, "\lsub{3}\C_5"'] & \R \arrow[l, "\lsub{5}\R_6"']
		\end{tikzcd}
	\end{equation*}
	\begin{equation*}
		\begin{tikzcd}[row sep = 0.3cm, column sep = 1cm]
			& 1 \arrow[r] \arrow[dd] & 2 \arrow[dd] & 3 \arrow[l] \arrow[dd] \\ 
			Q^1\tilde\otimes Q^2: \\
			& 4 \arrow[r] & 5 \arrow[luu] \arrow[ruu] & 6 \arrow[l]
		\end{tikzcd}
	\end{equation*}
	\begin{equation}\label{eq - potential in example at end 1}
		\begin{aligned}
			W(S^1, S^2) =& \lsub{2}1_1\otimes \lsub{1}1_5\otimes \lsub{5}1_2 - \lsub{2}1_1\otimes \lsub{1}i_5\otimes \lsub{5}i_2 + \lsub{2}1_3\otimes \lsub{3}1_5\otimes \lsub{5}1_2 - \lsub{2}1_3\otimes \lsub{3}i_5\otimes \lsub{5}i_2 + \\
			&- \lsub{4}1_1\otimes \lsub{1}1_5\otimes \lsub{5}1_4 - \lsub{6}1_3\otimes \lsub{3}1_5\otimes \lsub{5}1_6.
		\end{aligned}
	\end{equation}
	Due to Proposition~\ref{proposition - preprojective algebra of tensor species}, $\Pi_3(\Lambda) = \mathcal{P}(S(S^1, S^2), W(S^1, S^2))$ where $\Lambda = T(S^1)\otimes_\R T(S^2)$. Since $S^1$ and $S^2$ are both $2$-homogeneous representation finite species, and they are of type $A_3$ and $B_2$ respectively, we can apply Corollary~\ref{corollary - tensor product species morita transitive}. All $2$-APR tilts are given in Figure~\ref{figure - Ex 1}, where connections are given by cut-mutation.
	\begin{figure}[!ht]
	\begin{adjustbox}{width=\columnwidth,center}
		\begin{tikzcd}[bezier bounding box]
			& & \Image{ex1/1/1.pdf} \ar[dll, -, thick] \ar[-, thick, d] \ar[-, thick, drr]  \\
			\Image{ex1/2/2.pdf} \ar[-, thick, d] \ar[-, thick, dr] & & \Image{ex1/3/3.pdf} \ar[-, thick, dl] \ar[-, thick, d] \ar[-, thick, dr] & & \Image{ex1/4/4.pdf} \ar[-, thick, d] \ar[-, thick, dl] \\
			\Image{ex1/5/5.pdf} \ar[-, thick, d] & \Image{ex1/12/12.pdf} \ar[-, thick, dl] \ar[-, thick, dr] & \Image{ex1/11/11.pdf} & \Image{ex1/10/10.pdf} \ar[-, thick, dl] \ar[-, thick, dr] & \Image{ex1/7/7.pdf} \ar[-, thick, d] \\
			\Image{ex1/8/8.pdf} \ar[-, thick, from=urr, crossing over, out=210, in=0] & & \Image{ex1/13/13.pdf} & & \Image{ex1/9/9.pdf} \ar[-, thick, from=ull, crossing over, out=-30, in=180] \\
			& & \Image{ex1/6/6.pdf} \ar[-, thick, uuull, out=190, in=230, looseness=1.6] \ar[-, thick, ull] \ar[-, thick, uuurr, out=-10, in=-50, looseness=1.6] \ar[-, thick, urr] \ar[-, thick, u]
		\end{tikzcd}
	\end{adjustbox}
		\caption{Mutation lattice for $A_3\times B_2$}\label{figure - Ex 1}
	\end{figure}
\end{myex}

\begin{myex}
	Let $F = \K$ and $G$ be a finite dimensional division algebra over $\K$. Let $S^1$ be the species
	\begin{equation*}
		\begin{tikzcd}
			F \ar[r, "F"] & F \ar[r, "F"] & F \ar[d, "F"] & F \ar[l, "F", swap] & F \ar[l, "F", swap] \\
			& & F
		\end{tikzcd}
	\end{equation*}
	of type $E_6$ and $S^2$ be the species
	\begin{equation*}
		\begin{tikzcd}
			G \ar[r, "G"] & G \ar[r, "G"] & F \ar[r, "F"] & F
		\end{tikzcd}
	\end{equation*}
	of type $F_4$. Both $S^1$ and $S^2$ are $6$-homogeneous by Corollary~\ref{Corollary - l homogeneous species} and representation finite by Theorem~\ref{theorem - dlab ringel}. Hence $T(S^1)\otimes_\K T(S^2)$ is a $2$-representation finite $\K$-algebra. By applying the canonical isomorphisms $F\otimes_F F \cong F$ and $F\otimes_F G \cong G$ and applying Proposition~\ref{proposition - preprojective algebra of tensor species} we get $\Pi_3(T(S^1)\otimes_\K T(S^2))\cong \mathcal{P}(S, W)$. Since $M_\alpha = G$ if and only if $D_{s(\alpha)} = G$ or $D_{t(\alpha)} = G$ for $\alpha\in Q_1$ and $M_\alpha = F$ otherwise, we can omit writing out $M_\alpha$ in the diagram of $S$. The species $S$ is given as in Figure~\ref{figure - species e6 x f4} over the quiver in Figure~\ref{figure - quiver e6 x f4}.
	\begin{figure}
	\begin{minipage}[c]{0.45\linewidth}
		\includegraphics[width=\linewidth]{ex2/species/species.pdf}
	\caption{Species of type \\ $E_6\times F_4$}\label{figure - species e6 x f4}
	\end{minipage}
	\hfill
	\begin{minipage}[c]{0.45\linewidth}
		\includegraphics[width=\linewidth]{ex2/quiver/quiver.pdf}
	\caption{Quiver of type $E_6\times F_4$}\label{figure - quiver e6 x f4}
	\end{minipage}
	\end{figure}
	Using the notation that $\lsub{t(\alpha)}a_{s(\alpha)}\in M_\alpha$ for $\alpha\in Q_1$, the potential $W$ is
	\begin{equation*}
		\begin{aligned}
			W =& \sum_{\substack{i\in \{2, 4\} \\ j\in \{2, 3\}}} \left( \lsub{(i-1)j}1_{(i-1)(j-1)} \otimes \lsub{(i-1)(j-1)}1^*_{ij}\otimes \lsub{ij}1_{(i-1)j} - \lsub{i(j-1)}1_{(i-1)(j-1)} \otimes \lsub{(i-1)(j-1)}1^*_{ij}\otimes \lsub{ij}1_{i(j-1)} \right) + \\
			&+ \sum_{\substack{i\in \{2, 4\} \\ j\in \{3, 4\}}} \left( \lsub{(i-1)j}1_{(i-1)(j+1)} \otimes \lsub{(i-1)(j+1)}1^*_{ij}\otimes \lsub{ij}1_{(i-1)j} - \lsub{i(j+1)}1_{(i-1)(j+1)} \otimes \lsub{(i-1)(j+1)}1^*_{ij}\otimes \lsub{ij}1_{i(j+1)} \right) + \\
			&+\sum_{i\in \{2, 4\}} \left( \lsub{(i-1)6}1_{(i-1)3} \otimes \lsub{(i-1)3}1^*_{i6}\otimes \lsub{i6}1_{(i-1)6} - \lsub{i3}1_{(i-1)3} \otimes \lsub{(i-1)3}1_{i6}\otimes \lsub{i6}1_{i3} \right) + \\
			&+ \sum_{j\in \{2, 3\}} \left( \lsub{2j}1_{2(j-1)} \otimes \lsub{2(j-1)}1^*_{3j}\otimes \lsub{3j}1_{2j} + \lsub{2j}1_{2(j-1)} \otimes \lsub{2(j-1)}i^*_{3j}\otimes \lsub{3j}i_{2j} - \lsub{3(j-1)}1_{2(j-1)} \otimes \lsub{2(j-1)}1^*_{3j}\otimes \lsub{3j}1_{3(j-1)} \right) + \\
			&+ \sum_{j\in \{3, 4\}} \left( \lsub{2j}1_{2(j+1)} \otimes \lsub{2(j+1)}1^*_{3j}\otimes \lsub{3j}1_{2j} + \lsub{2j}1_{2(j+1)} \otimes \lsub{2(j+1)}i^*_{3j}\otimes \lsub{3j}i_{2j} - \lsub{3(j+1)}1_{2(j+1)} \otimes \lsub{2(j+1)}1^*_{3j}\otimes \lsub{3j}1_{3(j+1)} \right) + \\
			&+ \left( \lsub{26}1_{23} \otimes \lsub{23}1^*_{36}\otimes \lsub{36}1_{26} + \lsub{26}1_{23} \otimes \lsub{23}i^*_{36}\otimes \lsub{36}i_{26} - \lsub{33}1_{23} \otimes \lsub{23}1_{36}\otimes \lsub{36}1_{23} \right).
		\end{aligned}
	\end{equation*}
	To clarify, $Q_2$ consists of all $3$-cycles in $Q$. Choosing cuts yields $2$-representation finite $\K$-algebras. One may compute that there are a total of $16599$ of cuts for $(S, W)$. We can for example choose the cut $C$ which consist of all diagonal arrows, illustrated in Figure~\ref{figure - cut diagonal arrows}. Then the truncated Jacobian algebra $\mathcal{P}(S, W, C) \cong T(S^1)\otimes_\K T(S^2)$. Another example of a cut would be to pick all vertical arrows. It is straightforward to check that it is a cut, and therefore the truncated Jacobian algebra will be a $2$-representation finite $\K$-algebra. The cut is illustrated in Figure~\ref{figure - vertical arrow cut}.
	\begin{figure}
	\begin{minipage}[c]{0.45\linewidth}
		\includegraphics[width=\linewidth]{ex2/diagonalcut/dcut.pdf}
	\caption{Cut with all diagonal arrows}\label{figure - cut diagonal arrows}
	\end{minipage}
	\hfill
	\begin{minipage}[c]{0.45\linewidth}
		\includegraphics[width=\linewidth]{ex2/verticalcut/vcut.pdf}
	\caption{Cut with all vertical arrows}\label{figure - vertical arrow cut}
	\end{minipage}
	\end{figure}
\end{myex}

\iffalse
\section{2-Representation Finite algebras and Truncated Jacobian Algebras}\label{Section - 2-rf and jac alg}
\note{This section need to be rewritten for Kellers recipe in known cases}
Let $\Lambda$ be a $\K$-algebra such that $\mathrm{gl.dim}(\Lambda) \le 2$. We say that $\Lambda$ is $2$-representation finite if there exists a cluster tilting $\Lambda$-module $M$, i.e.
\begin{equation*}
	\begin{aligned}
		\mathrm{add} (M) &= \{X\in \mathrm{mod}\Lambda \mid \mathrm{Ext}^1_\Lambda(M, X) = 0 \} = \\
		&= \{X\in \mathrm{mod}\Lambda \mid \mathrm{Ext}^1_\Lambda(X, M) = 0 \}.
	\end{aligned}
\end{equation*}
It is shown in \cite{IO13Stablecateofhigherpreproj} that for a $2$-representation finite $\K$-algebra $\Lambda$ its $3$-preprojective algebra $\Pi_3(\Lambda)$ is a finite dimensional self-injective algebra.

\begin{myprop}\label{proposition - algebra species w relations}\cite[Proposition 3.3]{berg2011structure}
	Let $\K$ be a perfect field and let $\Lambda$ be a $\K$-algebra of global dimension at most $2$. Then $\Lambda \cong T(S)/R$ for some species $S$ and admissible relations $R$.
\end{myprop}

In the light of Proposition~\ref{proposition - algebra species w relations} we can assume that $\Lambda = T(S)/R$ where $S$ is some species and $R$ some admissible relations. The following theorem describes the relation between $\Pi_3(\Lambda)$ and species with potential.

\begin{mythm}\label{theorem - preproj as jacobian alg}
	Let $\Lambda = T(S)/R$ be a $\K$-algebra of global dimension at most $2$, where $S$ is a species over the quiver $Q$. We have an isomorphism $\Pi_3(\Lambda) \cong \mathcal{P}(S_\Lambda, W_\Lambda)$, where the species with potential $(S_\Lambda, W_\Lambda)$ is defined in the following steps:
	\begin{enumerate}
		\item The quiver $\tilde{Q}$ is defined as:
		\begin{enumerate}[label = -]
			\item $\tilde{Q}_0 = Q_0$.
			\item $\tilde{Q}_1 = Q_1 \cup \{i\to j \mid e_i R e_j \not= 0 \} = Q_1 \cup Q_R$.
		\end{enumerate}
		\item The species $S_\Lambda = (\tilde{D}, \tilde{M})$ over $\tilde{Q}$ is defined as:
		\begin{enumerate}[label = -]
			\item $\tilde{D} = D$.
			\item $\tilde{M} = M\oplus R^* = M\oplus \bigoplus_{\alpha\in Q_R} (e_{s(\alpha)Re_{t(\alpha)}})^*$
		\end{enumerate}
		\item The potential $W_\Lambda$ is defined as the image of $1_D$ under he counit
		\begin{equation*}
			D \to R\otimes_D R^*.
		\end{equation*}
	\end{enumerate}
\end{mythm}

\begin{myrem}
	The quiver case of Theorem~\ref{theorem - preproj as jacobian alg} is covered by \cite[Theorem 6.10]{Keller2011CYcompletions}.
\end{myrem}

\begin{proof}
	\note{I still don't have a proof for this. I still do not understand what happens with the differential on A x QV x A on page 14-15. Also need Lemma 2.8. Also if I only care about the homology at 0, can't I just defined the differential d however I want and then force the result?}
\end{proof}

We also need to following special case when our algebra $\Lambda$ is given by a truncated Jacobian algebra.

\begin{myprop}\label{proposition - S_A iso to S}
	Let $\Lambda = \mathcal{P}(S, W, C)$ for some species with potential $(S, W)$ and an algebraic cut $C$. Then $\Pi_3(\Lambda) \cong \mathcal{P}(S, W)$. Moreover, $(S, W)$ and $(S_\Lambda, W_\Lambda)$ are isomorphic as species with potentials.
\end{myprop}

\begin{proof}
	\textbf{extend prop before proving it? i.e. that $Q_R$ is a cut}
\end{proof}

\begin{mythm}
	Let $(S, W)$ be a self-injective species with potential and $C$ an algebraic cut. Then $\mathcal{P}(S, W, C)$ is a $2$-representation finite algebra.
\end{mythm}

\begin{proof}
	From the definition of algebraic cuts we have that $\Lambda = \mathcal{P}(S, W, C)$ has global dimension at most $2$. By \cite[Corollary 3.4]{IO13Stablecateofhigherpreproj} we have that $\Pi_3(\Lambda) = \mathcal{P}(S_\Lambda, W_\Lambda)$ is a finite dimensional self-injective species with potential if and only if $\Lambda$ is $2$-representation finite. The former follows from Proposition~\ref{proposition - S_A iso to S}.
\end{proof}
\fi

\clearpage
\bibliographystyle{alpha}
\bibliography{References}
\end{document}
